\documentclass[11pt,a4paper]{article}
\usepackage[utf8]{inputenc}
\usepackage[english]{babel}
\usepackage{amsmath}
\usepackage{amsfonts}
\usepackage{amssymb}
\usepackage{amsthm}
\usepackage{hyperref}
\usepackage{multirow}

\newenvironment{acknowledgements}{%
  \renewcommand{\abstractname}{Acknowledgements}% Rename Abstract to Acknowledgements
  \begin{abstract}
}{%
  \end{abstract}
}

%%%% an abstract-like acknowledgement environment
\newcommand{\ackname}{Acknowledgements}
\makeatletter
\if@titlepage
  \newenvironment{acknowledgement}{%
    \titlepage
    \null\vfil
    \@beginparpenalty\@lowpenalty
    \begin{center}%
      \bfseries \ackname
      \@endparpenalty\@M
    \end{center}}%
  {\par\vfil\null\endtitlepage}
\else
  \newenvironment{acknowledgement}{%
    \if@twocolumn
      \section*{\ackname}%
    \else
      \small
      \begin{center}%
        {\bfseries \ackname\vspace{-0.5em}\vspace{\z@}}%
      \end{center}%
      \quotation
    \fi}
    {\if@twocolumn\else\endquotation\fi}
\fi
\makeatother

% Theorem environments
\newtheorem{theorem}{Theorem}[section]
\newtheorem{corollary}[theorem]{Corollary}
\newtheorem{definition}[theorem]{Definition}
\newtheorem{lemma}[theorem]{Lemma}
\newtheorem{remark}[theorem]{Remark}

\allowdisplaybreaks

\title{Alpöge's counterexample to the JC, Economic Stabilization, and Zhao's Vanishing Conjecture}

\author{Michiel de Bondt}

\begin{document}

\maketitle

\begin{abstract}
In this paper, we apply stabilization techniques which are economic in the number of added dimensions to a slightly tweaked version of Alpöge's counterexample to $\mathrm{JC}(3)$ and $\mathrm{JC}(3,7)$. The first idea which lead to more economic stabilization is to bypass the degree 3 reduction of Bass, Connell and Wright in the construction of a counterexample to $\mathrm{HJC}(n)$. We construct counterexamples to $\mathrm{HJC}(15)$, $\mathrm{VC}(28)$, and $\mathrm{HVC}(30)$. On the other hand, we will perform degree 3 reduction after all to obtain counterexamples of low degree. We construct counterexamples to $\mathrm{JC}(12,3)$, $\mathrm{HJC}(18,3)$, $\mathrm{VC}(34,4)$, and $\mathrm{HVC}(36,4)$. As an encore, we construct a cubic-linear counterexample to the Jacobian Conjecture in dimension 90. All presented counterexamples to the Jacobian Conjecture (Vanishing Conjecture) are over the (Gaussian) integers.
\end{abstract}

% *************************************************************************************************

\section{Introduction}
\label{sec:introduction}

The Jacobian Conjecture, posed by Ott-Heinrich Keller in 1939, asserts that any polynomial map $F: \mathbb{C}^n \to \mathbb{C}^n$ with a nowhere-vanishing Jacobian determinant ($\det \mathcal{J}(F) \in \mathbb{C}^\times$) is globally invertible. We denote this classic conjecture in dimension $n$ as $\mathrm{JC}(n)$, and restrict it to polynomial maps of exact degree $d$ as $\mathrm{JC}(n,d)$. If the non-linear part of the mapping is strictly homogeneous of degree $d$, the corresponding assertion is denoted as $\mathrm{HJC}(n,d)$. 

For nearly nine decades, this problem stood as one of the most notorious challenges in algebraic geometry and polynomial algebra. In July 2026, Levent Alpöge published a historic breakthrough, constructing an explicit, non-injective polynomial counterexample establishing the falsity of $\mathrm{JC}(3)$ and, more specifically, $\mathrm{JC}(3,7)$ \cite{alpoge2026}.

Alpöge's counterexample immediately revitalized interest in several long-standing conjectures known to be stably-equivalent to the Jacobian Conjecture. Chief among these is the \emph{Zhao's Vanishing Conjecture} (for the Laplace operator), introduced by Wenhua Zhao in 2007 \cite{zhao2007}, which we will henceforth abbreviate as the \emph{Vanishing Conjecture}. The conjecture states that for the standard Laplace operator $\Delta = \sum \frac{\partial^2}{\partial x_k^2}$ and a polynomial $P \in \mathbb{C}[x_1,x_2,\ldots,x_n]$, if $\Delta^m(P^m) = 0$ for all $m \ge 1$, then $\Delta^m(P^{m+1}) = 0$ for all $m \gg 0$. We denote this assertion for polynomials in $n$ variables as $\mathrm{VC}(n)$, and specify the restriction to polynomials of exact degree $d$ as $\mathrm{VC}(n,d)$. When the underlying potential $P$ is strictly homogeneous (of degree $d$), the conjecture is denoted as $\mathrm{HVC}(n)$ ($\mathrm{HVC}(n,d)$). Zhao established that the universal (all-dimensional) Jacobian Conjecture is completely equivalent to the universal Vanishing Conjecture for (quartic) homogeneous polynomials under the standard Laplacian.

To systematically manipulate and stabilize these counterexamples, it is highly convenient to operate within the class of unipotent polynomial automorphisms. A polynomial map $F: \mathbb{C}^n \to \mathbb{C}^n$ is called a \emph{unipotent Jacobian map} if it takes the form $F = X + H$, where the Jacobian matrix $\mathcal{J}(H)$ is strictly nilpotent. We define $\mathrm{UJC}(n)$ as the statement that every unipotent Jacobian map over $\mathbb{C}$ in dimension $n$ is globally invertible, and denote its restriction to exact degree $d$ as $\mathrm{UJC}(n,d)$. By standard algebraic frameworks, any counterexample to $\mathrm{JC}(n,d)$ can be embedded into a higher-dimensional unipotent setting without altering its structural noninjectivity.

To transport a counterexample from the Jacobian framework to the vanishing arena, researchers traditionally rely on the classical Bass--Connell--Wright (BCW) reduction \cite{BCW1982}. The BCW procedure introduces a vast sequence of auxiliary variables to systematically force the degree of any Keller map down to 3, and homogenizing the non-linear part. In this paper, we proceed differently: we skip the intermediate degree-3 reduction entirely and synchronize the system by making the non-linear part homogeneous immediately, without altering the natural highest degree of the counterexample. As detailed in Section~\ref{sec:unipotent_reduction}, this direct, single-step homogenization preserves tight control over the variable expansion, terminating at an exact total dimension of:
\begin{equation} \label{eq:hmgM}
M = n + \sum_{i=1}^{n} \max(\deg(F_i) - 2, \, 0) + 1 = \mathcal{O}(n \cdot d)
\end{equation}
where $d = \max_{i} (\deg F_i)$ represents the maximum degree of the system. Applied to our volume-normalized version of Alpöge's mapping, this stabilization strategy yields explicit counterexamples to $\mathrm{UJC}(14)$, $\mathrm{UJC}(14,7)$, $\mathrm{HJC}(15)$, and $\mathrm{HJC}(15,7)$.

In contrast, forcing a degree-3 restriction does not possess such a direct, predictable complexity bound. Because the traditional BCW method requires splitting higher-degree components into extensive binary trees of lower-degree cofactors, the stable configuration space is widely expected to suffer from extreme dimensional inflation. Indeed, in a recent preprint by Long \cite{long2026}, in which the immediate aftermath of Alpöge's breakthrough is explored, a direct application of the cubic--homogeneous BCW reduction is performed without any optimization of the dimension outcome, leading to a stable homogeneous map in dimension 79. Based on these initial results, one might easily conclude that any degree-3 BCW reduction is inherently expensive and dimensionally prohibitive.

Heuristically, this assumption can be heavily nuanced by employing dimensionally economic stabilization techniques. As detailed in Section~\ref{sec:bcw_optimized}, by abandoning dynamic variable generation in favor of a fixed, pre-optimized initial grid and deploying a two-tier quotient sweeping phase based on the quadratic anchors $x_1 x_2$ and $x_1 x_3$, the algebraic overhead can be radically compressed. Through a careful term-by-term puzzle, the cubic--homogeneous reduction terminates at only 21 dimensions, built over a 14-dimensional cubic base map. This allows us to construct explicit low-degree counterexamples to $\mathrm{JC}(14,3)$, $\mathrm{UJC}(20,3)$, and $\mathrm{HJC}(21,3)$, revealing a surprising structural efficiency within the stabilized coordinate rings of the counterexample.

The dimensions $14$ and $21$ respectively seem astonishing compared to the corresponding numbers $39$ and $79$ by Long. But the impressiveness of these numbers vanished when I found out that Roy van Rijn had come to dimensions $12$ and $19$. Without studying Roy van Rijn's construction, I tried to improve my own construction. Because of Eq.~\eqref{eq:hmgM}, I knew that Roy van Rijn's $12$-dimensional construction had exactly $19 - 12 - 1 = 6$ cubic components, just as my $14$-dimensional construction ($21 - 14 - 1 = 6$). So I tried to find a $12$-dimensional cubic construction with $6$ cubic components, and I found one. It appeared to be essentially the same as Roy van Rijn's construction. A day later, I tried to find a $12$-dimensional cubic construction with only $5$ cubic components, since that would reduce to a cubic homogeneous construction in dimension $12 + 5 + 1 = 18$, and again, I found one. These constructions are also described in Section~\ref{sec:bcw_optimized}. It appeared that my last construction can by applied partially, to obtain counterexamples to $\mathrm{JC}(n,d)$, $\mathrm{UJC}(n,d)$ and $\mathrm{HJC}(n,d)$ for $d = 4, 5, 6$.

The Hessian Conjecture asserts that any Hessian-type polynomial map over $\mathbb{C}$ is globally invertible. We denote this conjecture in dimension $n$ as $\mathrm{HC}(n)$, and restrict it to polynomial maps of exact degree $d-1$ as $\mathrm{HC}(n,d)$. In the early 2000s, Guowu Meng \cite{meng2005} established a fundamental bridge between the Hessian Conjecture and the classic Jacobian Conjecture, proving that $\mathrm{HC}(2n)$ implies $\mathrm{JC}(n)$ and, more specifically, that $\mathrm{HC}(2n, d+1)$ implies $\mathrm{JC}(n, d)$. The bridge was also obtained in \cite{BE2005}, but here, we explicitly utilize Meng's bridge (see also \cite{B2014}). For instance, our counterexample to $\mathrm{JC}(12,3)$, can be systematically converted into a concrete counterexample to $\mathrm{HC}(24,4)$. This structural link was recently explored by Meng and Liang Yang \cite{mengyang2026}, who applied these original reduction methods directly to Alpöge's three-dimensional counterexample, thereby establishing the falsity of $\mathrm{HC}(6)$ and $\mathrm{HC}(6,8)$. Furthermore, in the same work, the authors successfully optimized their six-dimensional construction by exploiting the relatively simple underling structure of Alpöge's map, ultimately improving the boundary to achieve a highly compact counterexample to $\mathrm{HC}(5)$ and $\mathrm{HC}(5,14)$.

When these stable mappings are passed to the differential gradient domain via the Hessian reduction rules established in Section~\ref{sec:gradient_reduction}, they generate definitive counterexamples to the Vanishing Conjecture. Under the standard dimension-doubling embedding, our direct homogeneous method transforms the counterexample against $\mathrm{HJC}(15)$ into a uniform potential establishing the falsity of $\mathrm{HVC}(30)$. In parallel, our optimized cubic--homogeneous BCW path lifts the counterexample against $\mathrm{HJC}(18,3)$ to drive a uniform potential establishing the falsity of $\mathrm{HVC}(36,4)$.

Finally, it is crucial to observe that global homogeneity is by no means mathematically required to invalidate the Vanishing Conjecture; according to Zhao's general duality framework, a strict unipotent matrix structure is completely sufficient. By dropping the projectivized scale plane $t$, our direct stabilization track yields a non-homogeneous map establishing the falsity of $\mathrm{UJC}(14)$, which under the complexified gradient reduction transforms into a potential establishing the falsity of $\mathrm{VC}(28)$. In parallel, our optimized BCW path drives a non-homogeneous map establishing the falsity of $\mathrm{UJC}(17,3)$, which lifts to a potential establishing the falsity of $\mathrm{VC}(34,4)$.

Crucially, all counterexamples to the Jacobian Conjecture and Hessian Conjecture presented throughout this work are defined strictly over the integers, ensuring arithmetic and numerical purity. Furthermore, all presented counterexamples to the Vanishing Conjecture are strictly over the Gaussian integers.

As an encore, we demonstrate in section \ref{sec:druz_linearity} how our 18-dimensional cubic--homogeneous unipotent counterexample can be systematically funnelled into a strictly cubic--linear structure of Drużkowski type. For homogeneous counterexamples against $\mathrm{HJC}(n,d)$ where the terminal coordinate is isolated as $x_n$, a powerful classification theorem can be deployed to determine a feasible dimension for a corresponding stable power-linear mapping of degree $d$. Applying this theorem with $d=3$ directly to our structure allows us to chart a mathematically streamlined linearization pathway. However, while our previous term-by-term puzzles successfully mitigated variable inflation, the algebraic transition required to fully linearize the arguments of these cubic monomials still demands a massive expansion of the underlying coordinate framework. By executing this layered sequence of linear substitutions, the system is driven into a record-low, yet inherently expensive dimension of 90. This secondary reduction highlights that while intermediate unipotent rings can be tightly compressed, the final descent to pure Drużkowski linearity inevitably triggers a severe monomial expansion explosion.

Beyond the immediate structural breakthroughs, the exact numerical descent of our power-linear dimensions uncovered a series of remarkable historical and biographical parallels within the Jacobian framework. Prior to our awareness of Van Rijn's result, our initial baseline Dru\.{z}kowski linearization stood firmly at 123 variables. When the integration of our optimal 18-dimensional homogeneous map first compressed this boundary down to exactly 99 variables, it mirrored a famous numerical milestone from the doctoral work of Engelbert Hubbers. In his 1998 PhD thesis, Hubbers stabilized Pinchuk's  counterexample to the Strong Real Jacobian Conjecture to a cubic--linear mapping in dimension 2033, which he subsequently optimized in 1999 to exactly dimension 1999 (notice that any real-valued counterexample to the Jacobian Conjecture is also a counterexample to the Strong Real Jacobian Conjecture). Although this structural coincidence dissolved about a week later when our polarization refinements drove the bound down to 98, a second historical resonance manifested when subsequent algebraic pruning collapsed the variable space to exactly 91 dimensions. This number matched an April Fools' prank orchestrated by Arno van den Essen in 2006, who had sent an email to prominent researchers in the field claiming that the author had discovered a counterexample in dimension 91—owing to the fact that April 1st represents the 91st day of 2006 as a standard calendar year. Remarkably, even this final playful correspondence was bound to be overtaken by the algebraic reality of the counterexample: again about a week later, the boundary collapsed to its definitive minimum of 90, a historic threshold that mirrors the exact day of Van den Essen's original email transmission on March 30th, the 90th day of 2006.

\subsection{The Base Mapping and Volume-Preserving Transformation}

To initialize our reduction framework, we examine Alpöge's original three-dimensional counterexample $F: \mathbb{C}^3 \to \mathbb{C}^3$, defined as:
\begin{align}
F_1 &= (1+x_1 x_2)^3 x_3 + x_2^2(1+x_1 x_2)(4+3x_1 x_2) \\
F_2 &= x_2 + 3x_1(1+x_1 x_2)^2 x_3 + 3x_1 x_2^2(4+3x_1 x_2) \\
F_3 &= 2x_1 - 3x_1^2 x_2 - x_1^3 x_3
\end{align}
Expanding the determinant of its Jacobian yields $\det \mathcal{J}(F) = -2$. Thus, Alpöge's mapping multiplies the volume by $2$ and flips the orientation of the configuration space.

To bring this mapping into perfect harmony with standard unipotent reduction, we tweak Alpöge's map to construct a volume- and orientation-preserving variant whose linear part is exactly the identity matrix. We define our initial transformed Keller map $G: \mathbb{C}^3 \to \mathbb{C}^3$ analytically via the relation:
\begin{equation}
G(x_1, x_2, x_3) = \frac{1}{2} \Big(F_3(x_1, 2x_2, 2x_3), \, F_2(x_1, 2x_2, 2x_3), \, F_1(x_1, 2x_2, 2x_3)\Big)
\end{equation}
Evaluating this composition explicitly yields the three-dimensional base map for our reduction sequence:
\begin{align}
G_1 &= x_1 - 3x_1^2 x_2 - x_1^3 x_3 \\
G_2 &= 3v^2 x_1 x_3 + (9v^2 - 6v - 2)x_2 \\
G_3 &= v^3 x_3 + (6v^2 + 2v)x_2^2
\end{align}
where $v = 1 + 2x_1 x_2$. The structural validity and coordinate properties of this identity mapping have been fully verified.

\subsection{Numerical Collision Universality and Parameterized Loci}

Because the linear part of $G$ is the identity matrix, the geometric locus exhibits a non-injective failure where multiple preimages collapse into identical target zones. By evaluating the kernel split analytically, we isolate two distinct, parameterized collision paths running along the coordinate axes:

\begin{enumerate}
    \item \textbf{Collision Locus on the Third Axis:} Along the third coordinate axis, the target collision path can be parameterized as $(0, \, 0, \, \frac{t}{2})$. For any non-zero parameter $t$, this destination possesses three distinct preimages in the domain:
    \begin{equation}
    \!\!\!\!\!\!\!\left\{ \left(0, \, 0, \, \frac{t}{2}\right), \; \left(\frac{\sqrt{-t}}{2t}, \, \frac{3\sqrt{-t}}{2}, \, -13t\right), \; \left(-\frac{\sqrt{-t}}{2t}, \, -\frac{3\sqrt{-t}}{2}, \, -13t\right) \right\}
    \end{equation}
    Specializing this system by substituting $t = -\frac{1}{4}$ yields the target point $(0, \, 0, \, -\frac{1}{8})$ with the three distinct preimages:
    \begin{equation}
    \left\{ \left(0, \, 0, \, -\frac{1}{8}\right), \; \left(-1, \, \frac{3}{4}, \, \frac{13}{4}\right), \; \left(1, \, -\frac{3}{4}, \, \frac{13}{4}\right) \right\}
    \end{equation}
    which corresponds precisely to the transformation and normalization of Alpöge's original non-injective collision coordinates.
    
    \item \textbf{Collision Locus on the Second Axis:} Alternatively, the target path along the second coordinate axis is parameterized via $(0, \, \frac{t}{2}, \, 0)$. This family tracks a collapse where three distinct preimages map onto the single axis destination:
    \begin{equation}
    \left\{ \left(0, \, \frac{t}{2}, \, -2t^2\right), \; \left(-\frac{2}{t}, \, \frac{t}{2}, \, t^2\right), \; \left(\frac{2}{t}, \, -\frac{t}{4}, \, \frac{5t^2}{8}\right) \right\}
    \end{equation}
    Specializing this configuration by substituting $t = 2$ collapses the locus into the target destination $(0, \, 1, \, 0)$ and generates three distinct preimages, two of which are strictly integral:
    \begin{equation}
    \left\{ \left(0, \, 1, \, -8\right), \; \left(-1, \, 1, \, 4\right), \; \left(1, \, -\frac{1}{2}, \, \frac{5}{2}\right) \right\}
    \end{equation}
    which highlights the rich geometric layout of our base configuration.
\end{enumerate}

The above coordinate collisions are formally verified within the accompanying Maxima worksheet \texttt{section1.mac}, along with the corresponding collisions of Alpöge's original counterexample. The existence of these coordinate collisions establishes the absolute validity of the three-dimensional counterexample. In the following sections, we demonstrate how these non-injective loci are embedded and stabilized into higher-dimensional unipotent gradient rings, driving the explicit failure of the Vanishing Conjecture.

% *************************************************************************************************

\section{Reduction to a Unipotent Map}
\label{sec:unipotent_reduction}

In this section, we establish the algebraic framework required to transform an arbitrary counterexample of the three-dimensional Jacobian Conjecture into a higher-dimensional polynomial map whose Jacobian matrix is unipotent. We provide a structural theorem addressing both the non-homogeneous and homogeneous cases, and formalize the stabilization procedure.

\subsection{The Unipotent Reduction Theorem}

Let $R$ be a commutative ring, and let $X = (x_1, x_2, \dots, x_n)$ denote the standard coordinates in $R^n$. We consider polynomial maps $F: R^n \to R^n$ of the form $F = (F_1, F_2, \dots, F_n)$, where each $F_i \in R[X]$.

\begin{definition}[(Homogeneous) Unipotent Polynomial Map]
A polynomial map $F: R^n \to R^n$ is said to be a \emph{(homogeneous) unipotent polynomial map} if it can be written in the form:
\begin{equation}
F(X) = X + H(X)
\end{equation}
where $H(X) = (H_1, H_2, \dots, H_n)$ is a (homogeneous) polynomial vector field whose Jacobian matrix $\mathcal{J}(H)$ is strictly nilpotent. In the non-homogeneous setting, $H(X)$ may contain non-zero constant or linear terms. The identity matrix guarantees that $\mathcal{J}(F) = I_n + \mathcal{J}(H)$ remains unipotent.
\end{definition}

Traditional reduction techniques, such as the Bass-Connell-Wright (BCW) reduction, systematically force the degree of a Keller map down to $3$. However, this process might introduce a severe inflation of dummy variables, expanding the dimension exponentially. By relaxing the target degree bound and accepting higher-degree structures, we can introduce a controlled stabilization that preserves a significantly lower dimension.

\begin{theorem}[Unipotent Reduction and Exact Dimension Stabilization]
\label{thm:unipotent_main}
Let $F: R^n \to R^n$ be a polynomial map whose linear part is the identity map, and let $d$ be a target degree such that $d \ge \max(\deg(F),2)$. 
\begin{enumerate}
    \item \textbf{Non-Homogeneous Case:} There exists an explicit algebraic stabilization that constructs a polynomial map  $G: R^m \to R^m$ of degree $d$, where the exact stable dimension $m$ is given by:
    \begin{equation}
    m = n + \sum_{i=1}^{n} \max(\deg(F_i) - 2, 0)
    \end{equation}
    such that for $H = G - X$:
    \begin{equation}
    \det \big(\mathcal{J} ((X + t H)\big) = \big(\det (\mathcal{J} F)\big)\big|_{X = t X}
    \end{equation}
    where $X = (x_1, x_2, \ldots, x_m)$.
            
    Furthermore, for each $p \in R^n$, there exists $q \in R[t]^m$ of the form $(tp, *, \ldots, *)$, with coordinates $n+1, n+2, \ldots, m$ in $R\,t^2 \cup \cdots \cup R\,t^{d-1}$, such that
    \begin{equation}
    \big(X + t^{-1}H)(q) = (t\, F(p), 0, \ldots, 0)
    \end{equation}
    In particular, $q|_{t=1} = (p, *, \ldots, *)$ and $G(q|_{t=1}) = (F(p), 0, \ldots, 0)$.
    
    \item \textbf{Homogeneous Case:} There exists an explicit extension of this stabilization that constructs a strictly homogeneous unipotent map $G_h: R^{m+1} \to R^{m+1}$ of degree $d$ by introducing exactly one additional global homogenizing variable, elevating the exact dimension to $m+1$, such that for $H_h = G_h - X_h$:
     \begin{equation}
    \det \big(\mathcal{J} (X_h + t H_h)\big) = \big(\det (\mathcal{J} F)\big)\big|_{X_h = t x_{m+1}^{d-2} X_h}
    \end{equation}        
    and $(H_h)_{m+1} = 0$, where $X_h = (x_1, x_2, \ldots, x_m, x_{m+1})$.
    
    Furthermore, for each $p \in R^n$, there exists $q_h \in R[t]^{m + 1}$ of the form $(tp, *, \ldots, *, 1)$, with coordinates $n+1, n+2, \ldots, m$ in $R\,t^2 \cup \cdots \cup R\,t^{d-1}$, such that
    \begin{equation}
    (X_h + t^{-1}H_h)(q_h) = (t\, F(p), 0, \ldots, 0, 1)
    \end{equation}
    In particular, $q_h|_{t=1} = (p, *, \ldots, *, 1)$ and $G_h(q_h|_{t=1}) = (F(p), 0, \ldots, 0, 1)$.
    
\end{enumerate}
\end{theorem}

\begin{corollary}
\label{cor:unipotent_determinant}
If the polynomial map $F$ in Theorem~\ref{thm:unipotent_main} possesses a constant Jacobian determinant, then $\det \mathcal{J}(F) = 1$ and the stabilized mappings $G$ and $G_h$ are strictly unipotent Jacobian maps. That is, the Jacobian matrices $\mathcal{J}(H)$ and $\mathcal{J}(H_h)$ are strictly nilpotent.
\end{corollary}

\begin{proof}
As the linear part of $F$ is $X$, $\det \mathcal{J}(F) = 1$ follows. We first address the non-homogeneous mapping $G: R^m \to R^m$ with $H = G - X$. By substituting the condition $\det \mathcal{J}(F) = 1$ into the matrix relation of Theorem~\ref{thm:unipotent_main}, the characteristic polynomial of $\mathcal{J}(H)$ collapses via the identity:
\begin{align}
t^m &= t^m \cdot 1\big|_{t = -t^{-1}} = t^m \cdot \big(\det \mathcal{J}(X + tH)\big)\big|_{t=-t^{-1}} \nonumber \\
&= t^m \cdot \det \big(I_m - t^{-1}\mathcal{J}(H)\big) = \det \big(t I_m - \mathcal{J}(H)\big)
\end{align}
By the Cayley--Hamilton theorem, the vanishing of this characteristic polynomial over the spectrum confirms that $\mathcal{J}(H)$ is strictly nilpotent, and thus $G$ is unipotent.

For the homogeneous mapping $G_h: R^{m+1} \to R^{m+1}$ with $H_h = G_h - X_h$, an analogous evaluation holds by exploiting the algebraic properties of the extended determinant. Alternatively, since $G_h$ is strictly homogeneous of degree $d \ge 2$, substituting the core condition $\det \mathcal{J}(F) = 1$ under the grading profile yields the relation:
\begin{align}
t^{m+1} &= t^{m+1} \cdot 1\big|_{X_h = t^{-1/(d-1)} X_h} = t^{m+1} \cdot \big(\det \mathcal{J}(G_h)\big)\big|_{X_h = t^{-1/(d-1)} X_h} \nonumber \\
&= t^{m+1} \cdot \det \big(I_{m+1} - t^{-1}\mathcal{J}(H_h)\big) = \det \big(t I_{m+1} - \mathcal{J}(H_h)\big)
\end{align}
The Cayley--Hamilton theorem similarly implies that $\mathcal{J}(H_h)$ is strictly nilpotent, establishing that $G_h$ is a homogeneous unipotent mapping.
\end{proof}

\subsection{Construction of the stabilization}

To manipulate the degrees of these mappings dynamically without disrupting invertibility conditions, we embed the system into a polynomial ring extended by a homogenization parameter $t$. This parameter controls the scaling of the highest-degree homogeneous components during the iterative dimension-raising steps.

Because each successive reduction step introduces a new coordinate linked to a specific higher-degree term, the process proceeds without gaps, allowing for an exact determination of the stable dimension.

We construct the map $G$ of dimension $m$ iteratively by resolving components whose degrees exceed the targeted bound $d$. To establish a uniform grading profile, we embed the system into the extended ring $R[t][X]$ by evaluating the initial map under a $t$-deformed scale:
\begin{equation}
G(X) = \frac{1}{t^{d-2}}F(t^{d-2}X)
\end{equation}
Under this transformation, a homogeneous component of $F$ of degree $j$ acquires a total degree of $j(d-2) - (d-2) = (j-1)(d-2) + 1$ in $G$ for $j = 1, 2, \ldots, d$. This shows that while gaps exist between these allowed degree profiles, the structure of the grading is strictly preserved under the scaling operation.

Suppose a component $G_i$ has a maximum total degree $k = \deg(G_i) > d$. We extract its leading homogeneous part of degree $k$, denoted here as $G_{i}^{(k)}$. We introduce a new auxiliary variable $x_{m+1}$ and append the new coordinate relation $h$ defined as:
\begin{equation}
h = x_{m+1} + \frac{G_{i}^{(k)}}{t^{k-d}}
\end{equation}
Notice that because $G_{i}^{(k)}$ is homogeneous of degree $k$, the fraction possesses a total degree of $k - (k-d) = d$. The component $h$ therefore consists exclusively of terms of degree $1$ (from $x_{m+1}$) and degree $d$ (from the fraction), maintaining a single gap of size $d-1$ within the appended relation.

By appending $h$ to $G$ as a new coordinate condition, we can subtract the leading high-degree contribution from the original component via the replacement:
\begin{equation}
G_i \leftarrow G_i - h \cdot t^{k-d} = G_i - G_{i}^{(k)} - x_{m+1} t^{k-d}
\end{equation}
This substitution completely annihilates the leading homogeneous part $G_{i}^{(k)}$ of degree $k$. The remaining newly introduced term, $-x_{m+1} t^{k-d}$, has a total degree of exact size $1 + (k-d) = k - d + 1$. 

The maximum degree of this component has therefore decreased from $k$ to $k - d + 1$, which represents a strict reduction of exactly $k - (k - d + 1) = d - 1$. Furthermore, because $x_{m+1}$ is a fresh variable, it cannot interact with or be canceled by any pre-existing terms in $G_i$, guaranteeing that the degree drops by exactly $d-1$. 

Since each step reduces the target degree by this uniform invariant without generating irregular intermediate degree patterns, the iteration terminates precisely when all components satisfy the degree bound $d$. This gap-free reduction sequence yields the exact stable final dimension of $m = n + \sum \max(\deg(F_i) - 2, 0)$.

For the non-homogeneous mapping, setting $t=1$ specializes the system back to $R^m$. For the homogeneous mapping, substituting the parameter $t$ with the $(m+1)$-th coordinate $x_{m+1}$ synchronizes the degrees, generating a perfectly uniform homogeneous profile of degree $d$ across all $m+1$ components. 

Note that in the non-homogeneous setting, the resulting unipotent map may contain linear terms other than those of its identity linear part, as a consequence of stabilizing higher-degree cross-terms.

\subsection{Proof of Theorem~\ref{thm:unipotent_main}}
\label{subsec:proof_unipotent_main}

In this subsection, we establish the formal proof of the determinant identities and the preimage vector identities asserted in Theorem~\ref{thm:unipotent_main}. The reduction mechanism relies on a sequence of projective substitutions and algebraic matrix factorizations that bridge the stable configuration spaces.

\subsubsection{Proof of the Determinant Identities}

We begin by establishing the determinant relation for the homogeneous mapping $G_h: R^{m+1} \to R^{m+1}$ evaluated at the projectivized scale plane $x_{m+1} = 1$. Let $H_h = G_h - X_h$. By mapping the unipotent factorization matrix chains through the localized Jacobian configurations, the determinant collapses via the following chain of identities:
\begin{align*}
\det J(X + t^{d-1} H) 
&= \det J(X_h + t^{d-1} H_h(X, 1)) \\
&= \det J(X_h + t^{-1} H_h(t X, t)) \\
&= (\det J(t^{-1} X, x_{m+1}))|_{X_h = (t X, x_{m+1}) + H_h(tX, t)} \cdot {} \\
&\qquad (\det J(X_h + H_h(X, t)))|_{X_h = (t X, x_{m+1})} \cdot {} \\
&\qquad \det J(t X, x_{m+1}) \\
&= t^{-m} \cdot (\det J(t^{2-d} F (t^{d-2}X)))|_{X = t X} \cdot t^{m} \\
&= ((\det J F)|_{X = t^{d-2} X})|_{X = t X} \\
&= (\det J F)|_{X = t^{d-1} X}
\end{align*}
By executing the parameter substitution $t \leftarrow t^{1/(d-1)}$, this immediately establishes the fundamental non-homogeneous identity:
$$
\det \mathcal{J}(X + t H) = \big(\det \mathcal{J} F\big)\big|_{X = t X}
$$
In addition, the parameter substitution $t \leftarrow  t^{1/(d-1)} x_{m+1}$ guarantees that
\begin{equation} \label{eq:hmgdet}
\det \mathcal{J}\big(X_h + t x_{m+1}^{d-1} H_h(X, 1)\big) = \big(\det \mathcal{J} F\big)\big|_{X = t x_{m+1}^{d-1} X}
\end{equation}

To extend this identity to the full, unconstrained homogeneous coordinate frame $X_h = (x_1, \dots, x_m, x_{m+1})$, we exploit the algebraic grading profile of the homogeneous field $H_h(X_h)$, which possesses a uniform degree of $d$. By applying the Jacobian chain rule under projective rescaling, and using Eq.~\ref{eq:hmgdet}, the general homogeneous determinant factors as follows:
\begin{align*}
\det J(X_h + t H_h)
&= \det J(X_h + t x_{m+1}^d H_h(x_{m+1}^{-1} X, 1)) \\
&= (\det J(x_{m+1} X, x_{m+1}))|_{X_h = (x_{m+1}^{-1}X,x_{m+1}) + t x_{m+1}^{d-1} H_h(x_{m+1}^{-1} X, 1)} \cdot {} \\
&\qquad (\det J(X_h + t x_{m+1}^{d-1} H_h(X, 1)))|_{X_h = (x_{m+1}^{-1} X, x_{m+1})} \cdot {} \\
&\qquad \det J(x_{m+1}^{-1} X, x_{m+1}) \\
&= x_{m+1}^m \cdot ((\det JF)|_{X = t x_{m+1}^{d-1} X})|_{X = x_{m+1}^{-1} X} \cdot x_{m+1}^{-m} \\
&= (\det JF)|_{X = t x_{m+1}^{d-2} X}
\end{align*}
This completes the proof of the determinant scaling laws for both the non-homogenous and homogeneous tracks.

\subsubsection{Proof of the Preimage Vector Identities}

We now establish the explicit coordinate profiles for the stable vectors $q$ and $q_h$. Let the intermediate polynomial map and the parameterized coordinate vector tracking the stabilization procedure under the scaling variable $t$ be denoted as $G_t$ and $q_t$, respectively. The tracking framework is initialized at the base dimension via the relation:
\begin{equation}
G_t = t^{2-d} F\big(t^{d-2} X\big)
\end{equation}
This algebraic initialization dictates that the tracking vector must be initialized as $q_t = t^{2-d} p$, which guarantees the exact baseline evaluation:
\begin{equation}
G_t(q_t) = t^{2-d} F(p)
\end{equation}
Consequently, every individual coordinate entry of the rescaled vector $t^{-1} q_t$ is strictly contained within the localized monomial subring $R\, t^{1-d}$.

Suppose that during the iterative step of the stabilization algorithm, a new coordinate component $h$ possessing a $t$-degree of $d - i$ is appended to the system vector $G_t$. Since the component $h$ has an associated $X$-degree of $i$, evaluating this block under the initialized parameter layer yields a total $t$-degree of:
\begin{equation}
(d - i) + (2 - d) i = d + i - d i = 1 + (1 - d)(i - 1)
\end{equation}
The next stable coordinate entry $(q_t)_{m+1}$ is uniquely determined as the solution to the kernel equation $h(t^{2-d} p) = 0$ with respect to the dummy variable $x_{m+1}$. As a result of this factorization, the rescaled coordinate $t^{-1} (q_t)_{m+1}$ respects the filtration sequence:
\begin{equation}
t^{-1} (q_t)_{m+1} \in R (t^{1-d})^{i-1} \subseteq R (t^{1-d})^2 \cup R (t^{1-d})^3 \cup \cdots \cup R (t^{1-d})^{d-1}
\end{equation}
By executing this expansion recursively until the system reaches its final stable dimension $m$, the global evaluation vector satisfies the unipotent boundary condition:
\begin{equation}
t^{-1} G_t \big(t \cdot t^{-1} q_t\big) = t^{-1} G_t \big(q_t\big) = \big(t^{1-d} F(p), \, 0, \, 0, \, \dots, \, 0\big)
\end{equation}

To transport this result back to the target coordinate frame of Theorem~\ref{thm:unipotent_main}, we apply the parameter substitution $t \leftarrow t^{1/(1-d)}$ and define the localized vector $q = (t^{-1} q_t)\big|_{t = t^{1/(1-d)}}$. Under this map, the auxiliary coordinates $n+1, n+2, \dots, m$ of $q$ are strictly contained within the polynomial partition $R t^2 \cup \cdots \cup R t^{d-1}$. The non-homogeneous evaluation identity then maps out as:
\begin{align}
(X + t^{-1} H)(q) &= t^{1/(d-1)} G_{t^{1/(1-d)}}(t^{1/(1-d)} q) \nonumber \\
&= \big(t^{-1} G_t (t \cdot t^{-1} q_t)\big)\big|_{t = t^{1/(1-d)}} \nonumber \\
&= \big(t^{-1} t^{2-d} F(p), \, 0, \, 0, \, \dots, \, 0\big)\big|_{t = t^{1/(1-d)}} \nonumber \\
&= \big(t F(p), \, 0, \, 0, \, \dots, \, 0\big)
\end{align}

Finally, for the homogeneous track, we append the active scale parameter to the vector tail by setting $q_h = (q, \, 1)$. Utilizing the homogeneous degree-synchronization rules under $H_h$, the final extended projectivized identity reduces to:
\begin{equation}
(X_h + t^{-1} H_h) (q_h) = \big((X + t^{-1}H)(q), \; 1\big) = \big(t F(p), \, 0, \, 0, \, \dots, \, 0, \, 1\big)
\end{equation}
This completes the formal proof of the preimage coordinate behaviors for both stable dimensions.

\subsection{Algorithmic Correspondence and Base Configurations}

To demonstrate the concrete application of Theorem~\ref{thm:unipotent_main}, we present the coordinate configurations and loci for the three foundational mappings spanning dimensions 3, 14, and 15. The initial three-dimensional map is  the tweaked variant of Alpöge's counterexample from the introduction. We take the collision point with two integral pre-image points. The computations are performed in the ancillary Maxima worksheet \texttt{section2a.mac}.

\subsubsection{The Initial Three-Dimensional Mapping}
Let $F: \mathbb{C}^3 \to \mathbb{C}^3$ denote the tweaked three-dimensional Keller map of degree 7. This mapping is defined explicitly as:
$$
F = \big(x_1 - 3 x_1^2 x_2 - x_1^3 x_3, 3 v^2 x_1 x_3 + (9 v^2 - 6 v - 2) x_2,  v^3 x_3 + (6 v^2 + 2v) x_2^2\big)
$$
where $v = 1 + 2x_1x_2$. This mapping exhibits a non-injective collision where two distinct pre-images map to an identical point in the target space:
$$
F(0,1,-8) = (0,1,0) = F(-1,1,4)
$$

\subsubsection{The Stabilized Fourteen-Dimensional Unipotent Mapping}
By executing the non-homogeneous stabilization algorithm on $F_{(3)}$ with a target degree of $d=7$, we eliminate higher-degree cross-terms without gaps, yielding the stable unipotent polynomial map $G: \mathbb{C}^{14} \to \mathbb{C}^{14}$:
\begin{gather*}
G = \big(x_1-x_5, x_2-x_9+3x_1x_3, x_3-x_{14}+8x_2^2, \\
\qquad\quad x_4-x_1^3x_3, x_5-x_4-3x_1^2x_2, x_6+12x_1^3x_2^2x_3, \\
\qquad\quad x_7-x_6+36x_1^2x_2^3, x_8-x_7+12x_1^2x_2x_3, \\
\qquad\quad x_9-x_8+24x_1x_2^2, x_{10}+8x_1^3x_2^3x_3, \\
\qquad\quad x_{11}-x_{10}+24x_1^2x_2^4, x_{12}-x_{11}+12x_1^2x_2^2x_3, \\
\qquad\quad x_{13}-x_{12}+28x_1x_2^3, x_{14}-x_{13}+6x_1x_2x_3\big)
\end{gather*}
Under this transformation, the original three-dimensional configuration maps to the stable coordinates:
\begin{align*}
\lefteqn{G(0,1,-8,0,\ldots,0) = (0,1,0,0,\ldots,0)} \\
&\qquad= G(-1,1,4,-4,-1,48,12,-36,-12,32,8,-40,-12,12)
\end{align*}
and for $H = G - X$:
\begin{align*}
\lefteqn{(X + t^{-1}H)(0,t,-8t,0,\ldots,0) = (0,t,0,0,\ldots,0)} \\
&\qquad= (X + t^{-1}H)(-t,t,4t,-4t^3,-t^2,48t^5,12t^4,-36t^3,-12t^2, \\
&\qquad\hphantom{{}= (X + t^{-1}H)(-t,t,4t,{}}32t^6,8t^5,-40t^4,-12t^3,12t^2)
\end{align*}

\subsubsection{The Homogeneous Fifteen-Dimensional Unipotent Mapping}
Enforcing strict homogeneity via the introduction of the global scaling parameter $t = x_{15}$ yields a uniform degree 7 mapping $G_h: \mathbb{C}^{15} \to \mathbb{C}^{15}$:
\begin{gather*}
G_h = \big(x_1-x_5x_{15}^6, x_2-x_9x_{15}^6+3x_1x_3x_{15}^5, x_3-x_{14}x_{15}^6+8x_2^2x_{15}^5, \\
\qquad x_4-x_1^3x_3x_{15}^3, x_5-x_4x_{15}^6-3x_1^2x_2x_{15}^4, x_6+12x_1^3x_2^2x_3x_{15}, \\
\qquad x_7-x_6x_{15}^6+36x_1^2x_2^3x_{15}^2, x_8-x_7x_{15}^6+12x_1^2x_2x_3x_{15}^3, \\
\qquad x_9-x_8x_{15}^6+24x_1x_2^2x_{15}^4, x_{10}+8x_1^3x_2^3x_3, \\
\qquad x_{11}-x_{10}x_{15}^6+24x_1^2x_2^4x_{15}, x_{12}-x_{11}x_{15}^6+12x_1^2x_2^2x_3x_{15}^2,  \\
\qquad\quad~ x_{13}-x_{12}x_{15}^6+28x_1x_2^3x_{15}^3, x_{14}-x_{13}x_{15}^6+6x_1x_2x_3x_{15}^4, x_{15}\big)
\end{gather*}
The projective synchronization tracks the non-injective locus to:
\begin{align*}
\lefteqn{G_h(0,1,-8,0,\ldots,0,1) = (0,1,0,0,\ldots,0,1)} \\
&\qquad= G_h(-1,1,4,-4,-1,48,12,-36,-12,32,8,-40,-12,12,1)
\end{align*}
and for $H_h = G_h - X_h$:
\begin{align*}
\lefteqn{(X_h + t^{-1}H_h)(0,t,-8t,0,\ldots,0,1) = (0,t,0,0,\ldots,0,1)} \\
&\qquad= (X_h + t^{-1}H_h)(-t,t,4t,-4t^3,-t^2,48t^5,12t^4,-36t^3,-12t^2, \\
&\qquad\hphantom{{}= (X_h + t^{-1}H_h)(-t,t,4t,{}}32t^6,8t^5,-40t^4,-12t^3,12t^2,1)
\end{align*}


% *************************************************************************************************

\section{Reduction to Hessian-Type Maps and Zhao's Vanishing Conjecture}
\label{sec:gradient_reduction}

In this section, we transition from stable unipotent polynomial automorphisms to gradient maps of Hessian type. We first define a parameterized symmetric potential over the polynomial ring $\mathbb{C}[t]$ and evaluate its exact determinental profile. We then invoke Wenhua Zhao's duality theorem to establish that the core potential $h$ constitutes a definitive counterexample to the Vanishing Conjecture (VC). 

\subsection{The Parameterized Symmetric Potential and Determinant Valuation}
\label{subsec:symmetric_potential}

Let $S$ be a commutative ring with an element $i$ such that $i^2 + 1 = 0$. Let $F: S^n \to S^n$ be a stabilized unipotent polynomial map written as $F(X) = X + H(X)$, where $\mathcal{J}(H)$ is a nilpotent matrix. To analyze the behavioral deformation under linear scaling, we lift the system to a polynomial ring extended by a formal variable $t$, considering the parameterized mapping over $S[t][X]$ defined as:
\begin{equation}
F_t(X) = tX + H(X)
\end{equation}
Following the symmetric gradient reduction framework established by de Bondt and van den Essen \cite{BE2005}, we duplicate the dimension of the coordinate space by spanning a $2n$-dimensional affine structure with coordinates $(X,Y) = (x_1, \ldots, x_n, y_1, \ldots, y_n) \in S^{2n}$. We invoke the complexified independent linear frames $z_j = x_j + i y_j$ and their geoconjugate counters $\bar{z}_j = x_j - i y_j$. We define the structural symmetric potential $g_t(X,Y) \in S[t][X,Y]$ via the inner product pairing:
\begin{equation}
g_t(X,Y) := \sum_{j=1}^{n} \bar{z}_j \cdot (F_t)_j(z) = t \sum_{j=1}^{n} \bar{z}_j  z_j + \sum_{j=1}^{n} \bar{z}_j H_j(z) 
\end{equation}
which is of the form:
\begin{equation}
g_t(X,Y) = t \sum_{j=1}^{n} (x_j^2 + y_j^2) + h(X, Y) 
\end{equation}
where $h(X,Y) = \sum_{j=1}^n \bar{z}_j H_j(Z)$. Taking the real gradient of this full potential yields the parameterized Hessian mapping $\mathcal{G}_t = \nabla g_t: S^{2n} \to S^{2n}$:
\begin{equation}
\mathcal{G}_t(X,Y) = 2t(X,Y) + \nabla h(X,Y)
\end{equation}

By operating over the ring $S[t]$, we can determine the exact determinantal profile of the total Hessian matrix analytically, demonstrating that it respects the determinant condition of the Jacobian Conjecture for any unital specialization as long as $\frac12 \in S$.

\begin{theorem}[Determinant Invariance of the Symmetric Map]
\label{thm:hessian_determinant}
Let $g_t(X,\allowbreak Y) = t \sum (x_j^2 + y_j^2) + h(X,Y)$ be the parameterized potential generated by the deformed map $F_t = tX + H(X)$. Considered as a polynomial map over the ring $S[t]$, the determinant of the Hessian matrix $\mathcal{H}(g_t)$ is given exactly by the monomial:
\begin{equation}
\det \mathcal{H}(g_t) = (2t)^{2n}
\end{equation}
Consequently, if $\frac12 \in S$, then for any specialized value $t \in S^\times$, the Jacobian determinant of the mapping $\mathcal{G}_t(X,Y) = 2t(X,Y) + \nabla h(X,Y)$ is an invertible constant (which is the deteminant condition of the Jacobian conjecture).
\end{theorem}

\begin{proof}
We establish the exact relation between the real Hessian matrix $\mathcal{H}(g_t) = \mathcal{H}_{X,Y}(g_t)$ and the complexified Hessian matrix $\mathcal{H}_{Z,\bar{Z}}(g_t)$ by deploying the linear coordinate transformation defined by $z_j = x_j + i y_j$ and $\bar{z}_j = x_j - i y_j$. We formalize this transformation via the $2n \times 2n$ block matrix $T$:
\begin{equation}
T = 
\begin{pmatrix}
I_n & i I_n \\
I_n & -i I_n
\end{pmatrix}
\end{equation}
By computing the determinant of $T$ using standard block elimination, we find that $\det(T) = \det(-i I_n - i I_n) = \det(-2i I_n) = (-2i)^n$. Applying the multi-dimensional chain rule to the gradient vectors reveals that the real Hessian matrix satisfies the exact congruence relation:
\begin{equation}
\mathcal{H}_{X,Y}(g_t) = T^T \left( \mathcal{H}_{Z,\bar{Z}}(g_t)\big|_{Z = X + iY, \, \bar{Z} = X - iY} \right) T
\end{equation}
Notice that coordinate substitution does not alter the algebraic value of a constant determinant. So as soon as we infer that $\mathcal{H}_{Z,\bar{Z}}(g_t) \in \mathbb{C}[t]$ the determinant of the real Hessian matrix reduces strictly to:
\begin{equation}
\det \mathcal{H}_{X,Y}(g_t) = \det(T^T) \cdot \det \mathcal{H}_{Z,\bar{Z}}(g_t) \cdot \det(T) = \left[ \det(T) \right]^2 \cdot \det \mathcal{H}_{Z,\bar{Z}}(g_t)
\end{equation}
Evaluating the secondary partial derivatives of the potential $g_t = t \sum \bar{z}_j z_j + \sum \bar{z}_j H_j(Z)$ under the independent complex frames yields the block structure:
\begin{equation}
\mathcal{H}_{Z,\bar{Z}}(g_t) = 
\begin{pmatrix}
\sum_{m=1}^n \bar{z}_m \frac{\partial^2 H_m(Z)}{\partial z_k \partial z_j} & t I_n + \mathcal{J}(H)^T \\
t I_n + \mathcal{J}(H) & 0
\end{pmatrix}
\end{equation}
Because the bottom-right block is identically zero, the determinant collapses to the product of its off-diagonal blocks:
\begin{equation}
\det \mathcal{H}_{Z,\bar{Z}}(g_t) = (-1)^n \det\left( t I_n + \mathcal{J}(H)^T \right) \cdot \det\left( t I_n + \mathcal{J}(H) \right)
\end{equation}
Since transposition preserves characteristic structures and determinants, this simplifies via the nilpotency of $-\mathcal{J}(H)$. The characteristic polynomial yields $\det(t I_n + \mathcal{J}(H)) = t^n$, forcing:
\begin{equation}
\det \mathcal{H}_{Z,\bar{Z}}(g_t) = (-1)^n (t^n)^2 = (-1)^n t^{2n}
\end{equation}
We now combine the determinant of the transformation matrix $T$ with the complexified core result. Since $\left[ \det(T) \right]^2 = \left( (-2i)^n \right)^2 = (-4)^n$, the total real Hessian determinant evaluates to:
\begin{equation}
\det \mathcal{H}_{X,Y}(g_t) = (-4)^n \cdot (-1)^n t^{2n} = (4t^2)^n = (2t)^{2n}
\end{equation}
This confirms that for any specialization $t \in \mathbb{C}^\times$, the determinant is a non-vanishing constant, preserving the global Keller condition globally.
\end{proof}

\subsection{Connection between the Vanishing Conjecture and the Unipotent Jacobian Conjecture}
\label{subsec:vanishing_ujc_connection}

We begin by establishing the formal duality framework introduced by Wenhua Zhao \cite{zhao2007}, which connects the differential criteria of the Laplace operator to the unipotent inversion framework. Zhao established that if a polynomial potential $P \in \mathbb{C}[X]$ contains no terms of total degree $\le 1$, then $P$ serves as a counterexample to the Vanishing Conjecture if and only if the parameterized polynomial mapping $F_t(X) = X - t \nabla P$ constitutes a counterexample to the unipotent Jacobian Conjecture over the polynomial ring $\mathbb{C}[t]$. 

For our structural analysis, we observe that operating with the map $X - t \nabla P$ is algebraically equivalent to operating with $X + t \nabla P$ under a sign-inversion of the parameter space. To rigorously specialize this parameterized family into standalone numerical counterexamples over the base ring without risk of coordinate coalescence, one can deploy the following general theorem, which applies to polynomial mappings whose Jacobian matrices are not required to be symmetric.

\begin{theorem}
\label{thm:invertibility_specialization}
Let $R$ be a commutative ring, let $X = (x_1, \dots, x_n)$ denote the coordinate vector, and let $H: R^n \to R^n$ be a polynomial vector field. We consider the parameterized polynomial mapping $X + tH$ defined over the polynomial ring $R[t]$.
\begin{enumerate}

    \item If there exists a scalar element $\lambda \in R$ such that the specialized numerical mapping $X + \lambda H$ is not globally invertible over $R$, then the parameterized mapping $X + tH$ is not globally invertible over $R[t]$.
    
    \item If $\mathcal{J}(H)$ is strictly nilpotent, then the linear part of $X + tH$ is globally invertible over $R[t]$.
    
    \item If $R$ is a reduced ring where the zero polynomial is the unique polynomial over $R$ that evaluates identically to the zero function, and $X + tH$ is not globally invertible over $R[t]$ while the linear part of $X + tH$ is globally invertible over $R[t]$, then there exist infinitely many scalar values $\lambda \in R$ such that the specialized mapping $X + \lambda H$ is not globally invertible over $R$.
    
    \item If $X + tH$ is not globally invertible over $R[t]$ and the vector field $H$ is strictly homogeneous of degree $d \ge 2$, then for every element $\lambda \in R$ that is not a zero-divisor, the specialized mapping $X + \lambda H$ is not globally invertible over $R$.
    
\end{enumerate}
\end{theorem}

\begin{proof}
\begin{enumerate}

    \item We prove the assertion by contraposition. Suppose the parameterized mapping $X + tH$ is globally invertible over $R[t]$ with a well-defined polynomial inverse $G(X, t)$. Specializing the parameter value by substituting $t = \lambda$ yields the polynomial map $G(X, \lambda)$, which acts as the exact global inverse of the specialized map $X + \lambda H$ over $R$, contradicting the non-invertibility assumption.
    
    \item To establish the global invertibility of the linear part of $X + tH$ over $R[t]$, we examine its constant Jacobian configuration evaluated at the origin $X = 0$. The linear multiplier matrix of the parameterized system is given explicitly by $A = I_n + t \mathcal{J}(H)\vert_{X=0}$. Because the Jacobian matrix $\mathcal{J}(H)$ is strictly nilpotent by assumption, its evaluation $\mathcal{J}(H)\vert_{X=0}$ is likewise a nilpotent matrix over $R$, meaning there exists a minimal packing index $k \in \mathbb{N}$ such that $(\mathcal{J}(H)\vert_{X=0})^k = 0$. 
    
    Consequently, we can construct the formal inverse matrix via a finite Neumann series expansion:
    \begin{align}
    A^{-1} &= I_n - t \big(\mathcal{J}(H)\big|_{X=0}\big) + t^2 \big(\mathcal{J}(H)\big|_{X=0}\big)^2 - {} \nonumber \\
    &\qquad t^3 \big(\mathcal{J}(H)\big|_{X=0}\big)^3 + \dots \pm t^{k-1} \big(\mathcal{J}(H)\big|_{X=0}\big)^{k-1}
    \end{align}
    Since this matrix summation terminates at a finite degree $t^{k-1}$, all its entries are strictly polynomial components contained within the ring $R[t]$. Evaluating the product verifies that $A \cdot A^{-1} = I_n$. Because the linear part possesses a well-defined, polynomial inverse matrix over $R[t]$, it is globally invertible over the domain $R[t]$, completing the proof.
    
    \item Just as above, $A = I_n + t \mathcal{J}(H)\vert_{X=0}$ denotes the matrix representing the linear part of the parameterized mapping $X + tH$. Since $A$ is invertible, the exact inverse of $A$ over $R$ exists, and we denote it by $A^{-1}$. We define the modified polynomial mapping 
    \[ \Phi(X, t) = A\cdot\big(X + t(H-H(0))\big)
    \]
    By construction, the Jacobian matrix of this composition satisfies $\mathcal{J}(\Phi)\vert_{X=0} = A^{-1} A = I_n$, meaning the constant part of the Jacobian determinant is normalized to the identity matrix. Since \(X + tH\) is not globally invertible, the composition $\Phi(X, t)$ is likewise non-invertible. 
    
    According to the formal inversion criterion for polynomial automorphisms, the formal power series inverse of $\Phi(X, t)$ must contain at least one non-vanishing higher-order term whose total degree strictly exceeds the standard geometric bound $d^{n-1}$, where $d = \max_i(\deg \Phi_i)$. The coefficient of this higher-order term can be isolated as a non-zero polynomial $f(t) \in R[t]$. Since $R$ respects the functional identity property where the zero polynomial is the unique representative of the zero function, there must exist a scalar value $\lambda \in R$ such that $f(\lambda) \neq 0$. This implies that the formal power series inverse of the specialized mapping $\Phi(X, \lambda)$ retains a non-vanishing term beyond the geometric degree bound, proving that $X + \lambda H$ is not globally invertible over $R$. Furthermore, since the modified polynomial $(t - \lambda)f(t) \in R[t]$ is non-zero, there must exist a distinct scalar value $\lambda' \neq \lambda$ such that $f(\lambda') \neq 0$, generating a distinct non-invertible specialized map. Repeating this construction inductively establishes the existence of infinitely many non-invertible scalar specializations over $R$.
    
        \item We prove the assertion by contraposition. Suppose there exists an element $\lambda \in R$ that is not a zero-divisor such that the specialized numerical mapping $X + \lambda H$ is globally invertible over $R$. Because the vector field $H$ is strictly homogeneous of degree $d \ge 2$, applying the coordinate rescaling $X \leftarrow tX$ yields the homogeneity relation $H(tX) = t^d H(X)$. It follows that the projectivized map $t^{-1}(tX + \lambda H(tX)) = X + t^{d-1}\lambda H(X)$ is globally invertible over the polynomial ring $R[t^{d-1}\lambda]$. Consequently, $X + t\lambda H(X)$ is globally invertible over the polynomial ring $R[t\lambda]$.
    
    To lift this property to the full polynomial ring $R[t]$, we construct the multiplicative set $\Lambda = \{\lambda, \lambda^2, \lambda^3, \dots\}$. Since $\lambda$ is not a zero-divisor, $\Lambda$ contains no zero-divisors, making the canonical localization homomorphism from $R$ into the ring of fractions $\Lambda^{-1}R$ strictly injective. We may therefore naturally view $R[t]$ as a subring of the localized polynomial ring $(\Lambda^{-1}R)[t]$. Over this localized domain, the element $\lambda$ is a unit. We can then define the localized ring endomorphism $\phi: (\Lambda^{-1}R)[t] \to (\Lambda^{-1}R)[t]$ determined by the mapping $f(t) \mapsto f(\lambda t)$. Because $\lambda$ is invertible in $\Lambda^{-1}R$, this endomorphism is a strict automorphism possessing the exact algebraic inverse $\phi^{-1}: f(t) \mapsto f(\lambda^{-1} t)$. The restriction of $\phi$ to $R[t]$ is therefore a strict isomorphism between $R[t]$ and $R[t\lambda]$. This forces the core parameterized mapping $X + tH$ to be globally invertible over the standard polynomial ring $R[t]$, completing the proof. \qedhere
    
\end{enumerate}
\end{proof}

\subsection{Transform of the Collision}
\label{subsec:transform_collision}

In this subsection, we formalize the algebraic transformation of the non-injective locus under the complexified gradient embedding. To preserve the arithmetic purity of the coordinates, we track the geometric collapse directly through the parameterized configurations derived from our unipotent stabilization.

The baseline coordinate collision point for the parameterized mapping $X + t^{-1}H$ evaluates exactly to the vector:
\begin{equation}
(0, \, t, \, 0, \, 0, \, 0, \, 0, \, 0, \, 0, \, 0, \, 0, \, 0, \, 0, \, 0, \, 0) \in S^{14}
\end{equation}
in the non-homogeneous track, and expands to the projectivized scale plane as:
\begin{equation}
(0, \, t, \, 0, \, 0, \, 0, \, 0, \, 0, \, 0, \, 0, \, 0, \, 0, \, 0, \, 0, \, 0, \, 1) \in S^{15}
\end{equation}
in the homogeneous track. By executing the reciprocal parameter scaling, the corresponding target collision point for the perturbed system $tX + H$ shifts to:
\begin{equation}
(0, \, t^2, \, 0, \, 0, \, 0, \, 0, \, 0, \, 0, \, 0, \, 0, \, 0, \, 0, \, 0, \, 0) \in S^{14}
\end{equation}
in the non-homogeneous case, and synchronizes cross the grading frame to:
\begin{equation}
(0, \, t^2, \, 0, \, 0, \, 0, \, 0, \, 0, \, 0, \, 0, \, 0, \, 0, \, 0, \, 0, \, 0, \, t) \in S^{15}
\end{equation}
in the homogeneous case. Immersing these targets into the differential gradient domain requires appending exactly 14 and 15 zero-components on the front, respectively, and multiplying the resulting configurations by the transposed embedding matrix $T^T$. This algebraic transformation maps the numerical target collision locations in the gradient ring explicitly onto:
\begin{equation}
\begin{array}{l}
(0, \, \hphantom{-}t^2, \, \hphantom{i}0, \, 0, \, 0, \, 0, \, 0, \, 0, \, 0, \, 0, \, 0, \, 0, \, 0, \, 0, \\
\hphantom{(}0, \, -i t^2, \, 0, \, 0, \, 0, \, 0, \, 0, \, 0, \, 0, \, 0, \, 0, \, 0, \, 0, \, 0)
\end{array} \in S^{28}
\end{equation}
for the non-homogeneous system, and evaluates to:
\begin{equation}
\begin{array}{l}
(0, \, \hphantom{-}t^2, \, \hphantom{i}0, \, 0, \, 0, \, 0, \, 0, \, 0, \, 0, \, 0, \, 0, \, 0, \, 0, \, 0, \, \hphantom{-}t, \\
\hphantom{(}0, \, -i t^2, \, 0, \, 0, \, 0, \, 0, \, 0, \, 0, \, 0, \, 0, \, 0, \, 0, \, 0, \, 0, \, -i t)
\end{array} \in S^{30}
\end{equation}
for the homogeneous system.

To map the corresponding preimage trajectories into the complexified coordinate frame, let $p$ denote a valid preimage vector of the unipotent tracking locus mapping onto the base collision under $X + t^{-1}H$. In both the non-homogeneous and homogeneous tracks, the stable complex gradient preimage vector $q$ is uniquely determined by solving the system of coupled linear constraints $z_j = p_j$ and $\bar{z}_j = 0$. The vanishing of the complex conjugate branch, $\bar{z}_j = 0$, dictates the core relation $q_j = i q_{j+n}$. Substituting this dependency directly into the primary constraint plane yields the system of equations:
\begin{equation}
2 q_j = q_j + i q_{j+n} = p_j
\end{equation}
Isolating the independent coordinate fields maps out the exact complex preimage transformation rule:
\begin{equation}
q_j = \frac{1}{2} p_j \quad \text{and} \quad q_{j+n} = -\frac{1}{2} i p_j
\end{equation}
The non-injective structural collapses of this complete collision transformation are formally verified within the accompanying ancillary script file \texttt{section3a.mac}.

Finally, this parameterized structure reveals a fundamental arithmetic trade-off inherent to the reduction ring. If we specialize the parameter space by substituting $t = 1$, the construction yields a valid polynomial mapping defined strictly over the Gaussian integers $\mathbb{Z}[i]$ that represents the exact gradient of a scalar potential over $\mathbb{Z}[i]$; however, under this choice, the map is not unipotent. Conversely, if we select the specialization $t = \frac{1}{2}$, the system settles into a strictly unipotent polynomial automorphism over $\mathbb{Z}[i]$, but its components lose their gradient-type symmetry over the ring $\mathbb{Z}[i]$. We can repair the underlying scalar potential by replacing the base ring $\mathbb{Z}[i]$ by $\mathbb{Z}[i][\frac12]$.

% *************************************************************************************************

\section{Optimized Bass--Connell--Wright Reduction to Degree 3}
\label{sec:bcw_optimized}

In this section, we formalize an optimized version of the Bass--Connell--Wright (BCW) degree-reduction algorithm \cite{BCW1982}. We present our reduction framework as an evolving multi-stage refinement. First, we outline our original reduction to degree 3 without homogeniety, which is a more straightforward baseline reduction sequence that establishes a clear, accessible path to degree 3, terminating at a stable cubic map of 14 variables. This reduction is executed in \texttt{section4a.mac}. While not achieving a record-low variable count, this 14-dimensional cubic framework is mathematically transparent and serves as the conceptual foundation. 

In the subsequent subsections, we demonstrate how deeper insights into structural cofactors allow us to compress this space to 12 dimensions, by providing two additional reductions. The first reduction is essentially the same as that of Van Rijn. But the second reduction differs in that the subsequent reduction to cubic-homogeneous form, as outlined in section 2, ends in dimension 18 instead of 19. 

\subsection{Base Extension and Initial Monomial Grid}

Let $(G_1, G_2, G_3)$ be the volume-normalized three-dimensional counterexample to the Jacobian Conjecture defined in Section~\ref{sec:introduction}. To initiate the reduction sequence, we expand the affine coordinate space by appending exactly seven baseline auxiliary components representing six monomials in $x_1, x_2, x_3$, and the monomial $x_2 x_6$. The core monomials are:
\begin{equation}
x_1^2, \, x_1 x_2, \, x_1 x_3, \, x_2 x_6, \;\; x_1^2 x_3, \, x_1 x_2^2, \;\; x_1^2 x_2 x_3
\end{equation}
By ordering these components strictly from highest to lowest total degree, we define the initial extended mapping vector:
\begin{align}
G_4  &= x_4  + x_1^2 x_2 x_3 \\
G_5  &= x_5  + x_1^2 x_3 \\
G_6  &= x_6  + x_1 x_2^2 \\
G_7  &= x_7  + x_1^2 \\
G_8  &= x_8  + x_1 x_2 \\
G_9  &= x_9  + x_1 x_3 \\
G_{10} &= x_{10} + x_2 x_6
\end{align}
This forms our initial 10-dimensional extended polynomial mapping vector $G = (G_1, \dots, G_{10})$. The inclusion of the cross-interaction link $G_{10} = x_{10} + x_2 x_6$ ensures that the network remains algebraically closed under localized factorizations.

\subsection{The Monomial Factorization Loop and Early Termination}

The reduction algorithm operates component-by-component over the first four dimensions of $G$. For each component $G_c$ ($1 \le c \le 4$), the leading non-linear monomial $M$ (with total degree $d$ and scalar coefficient $c_M$) is captured sequentially. To minimize variable expansion, the core factorization loop differentiates between the components:
\begin{itemize}
    \item For $G_1$ and $G_4$, the loop runs iteratively until the component is reduced to \textbf{degree 3}.
    \item For $G_2$ and $G_3$, the loop terminates early, leaving the components at \textbf{degree 4}.
\end{itemize}

During this loop, the selected monomial $M$ is evaluated to find a divisor within the initial grid of degree $\lceil d/2 \rceil$. Because the grid is pre-optimized, the routine always locates a complete pair of indices $[i, j]$ from the baseline set ($1 \le i, j \le 7$) such that the product of the corresponding core monomials matches $M$ exactly. The component is then immediately updated via a direct factorization substitution:
\begin{equation}
G_c \leftarrow G_c - c_M \cdot G_{i+3} \cdot G_{j+3}
\end{equation}
where $G_{i+3}$ and $G_{j+3}$ represent the corresponding shifted variables from the baseline grid. Because the maximum degree of $G_4$ is $4$, its target ceiling bound is $\lceil 4/2 \rceil = 2$. Consequently, $G_4$ will only look for quadratic grid links and can never select its own core monomial as a divisor, rendering circular definitions impossible.

\subsection{The Secondary Quotient Sweeping Step for $G_2$ and $G_3$}

Because the primary factorization loop for $G_2$ and $G_3$ terminates at degree 4, these components require a dedicated secondary stabilization step to safely clear their remaining degree-4 excesses down to degree 3. This is achieved via exact algebraic polynomial division using the specific quadratic baseline anchors $x_1 x_2$ and $x_1 x_3$:

\begin{enumerate}
    \item \textbf{The \mathversion{bold}$x_1 x_2$-Quotient Sweep:} The algorithm isolates the homogeneous degree-4 part of the active component and divides it by $x_1 x_2$. The resulting quotient (which is strictly of degree 2) is bound to a new auxiliary variable by incrementing the stable dimension $m \leftarrow m + 1$:
    \begin{equation}
    G_m = x_m + Q_1
    \end{equation}
    where $Q_1$ is the algebraic quotient. The component is then updated by subtracting the product of this new variable and the $G_8 = x_8 + x_1 x_2$ grid anchor:
    \begin{equation}
    G_c \leftarrow G_c - G_8 \cdot G_m
    \end{equation}
    
    \item \textbf{The \mathversion{bold}$x_1 x_3$-Quotient Sweep:} Any remaining remainder from the previous division is subsequently divided by $x_1 x_3$. The new degree-2 quotient triggers another dimension increment $m \leftarrow m + 1$:
    \begin{equation}
    G_m = x_m + Q_2
    \end{equation}
    where $Q_2$ is the secondary quotient. The remaining degree-4 profile is completely annihilated by subtracting the product with the $G_9 = x_9 + x_1 x_3$ grid anchor:
    \begin{equation}
    G_c \leftarrow G_c - G_9 \cdot G_m
    \end{equation}
\end{enumerate}
Since both quotients are strictly quadratic, their multiplication with the quadratic baseline anchors generates pure degree-4 contributions that perfectly eliminate the highest-degree terms without introducing higher-order cross-contaminations, successfully forcing $G_2$ and $G_3$ to degree 3.


\subsection{The Resulting Fourteen-Dimensional Configuration}
\label{subsec:cubic_14d_baseline}

Following the systematic execution of our baseline reduction algorithm, the primary loop successfully clears the high-degree excesses of the counterexample by extending the coordinate space to exactly 14 dimensions. Because all appended auxiliary components evaluate identically to zero under the initialization kernel, the original non-injective locus is preserved perfectly. The complete algebraic system profiles, transformation mappings, and stable integer preimages are duplicated textually from the symbolic verification environment as follows:

\begin{itemize}
    \item \textbf{The Initial Extended Map \mathversion{bold}$F$:}
    \begin{align*}
    F = \big( & x_1 - 3x_1^2x_2 - x_1^3x_3, \; 3v^2x_1x_3 + (9v^2-6v-2)x_2, \\ 
    & v^3x_3 + (6v^2+2v)x_2^2, \; x_4 + x_1^2x_2x_3, \; x_5 + x_1^2x_3, \; x_6 + x_1x_2^2, \\ 
    & x_7 + x_1^2, \; x_8 + x_1x_2, \; x_9 + x_1x_3, \; x_{10} + x_2x_6, \\
    & x_{11} + 12x_1x_3 - 12x_2x_5 - 36x_2x_8, \; x_{12} - 12x_1x_6, \\
    & x_{13} + 28x_2^2 - 8x_2x_4 - 48x_2x_6 - 12x_2x_9, \; x_{14} + 8x_1x_{10} \big)
    \end{align*}
    
    \item \textbf{The Stable Geometric Preimages \mathversion{bold}$P$:}
    \begin{align*}
    P = \big\{ & (0, \, 1, \, -8, \, 0, \, 0, \, 0, \, 0, \, 0, \, 0, \, 0, \, 0, \, 0, \, -28, \, 0), \\
               & (-1, \, 1, \, 4, \, -4, \, -4, \, 1, \, -1, \, 1, \, 4, \, -1, \, 36, \, -12, \, 36, \, -8) \big\}
    \end{align*}
    
    \item \textbf{The Reduction Duality Identity Vector \mathversion{bold}$E$:}
    \begin{align*}
    E = \big( & y_1 + y_7y_9, \; y_2 - 12y_5y_6 - 36y_6y_8 - y_8y_{11} - y_9y_{12}, \\
              & y_3 - 8y_4y_6 - 24y_6^2 - 12y_6y_9 + 8y_5y_{10} - y_8y_{13} - y_9y_{14}, \\
              & y_4 - y_8y_9, \; y_5, \; y_6, \; y_7, \; y_8, \; y_9, \; y_{10}, \; y_{11}, \; y_{12}, \; y_{13}, \; y_{14} \big)
    \end{align*}
    
    \item \textbf{The Resulting Cubic Map \mathversion{bold}$F_{\text{cub}} = E(F)$:}
    \begin{align*}
    F_{\text{cub}} = \big( & x_1 - 3x_1^2x_2 + x_7x_9 + x_1^2x_9 + x_1x_3x_7, \\
    & x_2 + 3x_1x_3 + 24x_1x_2^2 - 12x_5x_6 - 36x_6x_8 - 36x_1x_2x_6 \\
    & \qquad - x_8x_{11} - x_1x_2x_{11} - 12x_1x_3x_8 + 12x_2x_5x_8 \\
    & \qquad + 36x_2x_8^2 - x_9x_{12} - x_1x_3x_{12} + 12x_1x_6x_9, \\
    & x_3 + 8x_2^2 + 6x_1x_2x_3 - 8x_4x_6 - 24x_6^2 - 12x_6x_9 \\
    & \qquad - 12x_1x_3x_6 + 8x_5x_{10} + 8 x_2 x_5 x_6 - x_8x_{13} \\
    & \qquad - x_1x_2x_{13} - 28x_2^2x_8 + 8x_2x_4x_8 + 48x_2x_6x_8 \\
    & \qquad + 12x_2x_8x_9 - x_9x_{14} - x_1x_3x_{14} - 8x_1x_9x_{10}, \\
    & x_4 - x_8x_9 - x_1x_2x_9 - x_1x_3x_8, \; x_5 + x_1^2x_3, \; x_6 + x_1x_2^2, \\
    & x_7 + x_1^2, \; x_8 + x_1x_2, \; x_9 + x_1x_3, \; x_{10} + x_2x_6, \\
    & x_{11} + 12x_1x_3 - 12x_2x_5 - 36x_2x_8, \; x_{12} - 12x_1x_6, \\
    & x_{13} + 28x_2^2 - 8x_2x_4 - 48x_2x_6 - 12x_2x_9, \; x_{14} + 8x_1x_{10} \big)
    \end{align*}
\end{itemize}

\subsection{A Twelve-Dimensional Configuration with Six Cubic Components}
\label{subsec:cubic_12d_six}

Following reports of a cubic reduction achieving low dimensions in contemporary literature, we systematically improved our cubic construction by introducing a highly compact baseline grid, resulting in a 12-dimensional coordinate system with six cubic components. This structure turns out to be essentially equivalent to the construction discovered by Roy van Rijn.

An inspection of our initial 14-dimensional framework reveals that the final reduction step executes a division by the anchor $x_1 x_3$ at a stage where only a single term of degree 4 remains in both $G_3$ and $G_4$. By expanding the initial baseline grid from seven to nine core monomials, and combining the new core monomials linearly with existing ones, we can actively prevent these specific high-degree residuals from manifesting at the end of the loop, thereby eliminating the need for the secondary quotient sweeping step and saving exactly two dimensions. The complete system parameters, evaluation mappings, and stable integer preimages are extracted directly from our symbolic worksheet \texttt{section4b.mac} as follows:

\begin{itemize}
    \item \textbf{The Initial Extended Map \mathversion{bold}$F$:} 
    \begin{align*}
    F = \big( & x_1 - 3x_1^2x_2 - x_1^3x_3, \; 3v^2x_1x_3 + (9v^2-6v-2)x_2, \\ 
    & v^3x_3 + (6v^2+2v)x_2^2, \; x_4 + x_1^2x_2^2, \; x_5 + x_1^2x_2, \; x_6 + x_1x_2x_3 + 3x_2^2, \\
    & x_7 + x_1^2, \; x_8 + x_1x_2, \; x_9 + x_1x_3, \; x_{10} + 3x_2x_3 - 2x_2 x_6, \\
    & x_{11} + x_1x_3 - x_1x_6 - x_3x_5, \; x_{12} + 7x_2^2 - 2x_3x_4 - x_1x_{10} \big)
    \end{align*}
    where $v = 1 + 2x_1x_2$.
    
    \item \textbf{The Stable Geometric Preimages \mathversion{bold}$P$:}
    \begin{align*}
    P = \big\{ & (0, \, 1, \, -8, \, 0, \, 0, \, -3, \, 0, \, 0, \, 0, \, 18, \, 0, \, -7), \\
               & (-1, \, 1, \, 4, \, -1, \, -1, \, 1, \, -1, \, 1, \, 4, \, -10, \, -1, \, -5) \big\}
    \end{align*}
    
    \item \textbf{The Reduction Duality Identity Vector \mathversion{bold}$E$:}
    \begin{align*}
    E = \big( & y_1 + y_7y_9, \; y_2 - 12y_5y_6 - 12y_8y_{11}, \; y_3 - 8y_4y_6 - 4y_5y_{10} - 4y_8y_{12}, \\
              & y_4 - y_8^2, \; y_5, \; y_6, \; y_7, \; y_8, \; y_9, \; y_{10}, \; y_{11}, \; y_{12} \big)
    \end{align*}
    
    \item \textbf{The Resulting Cubic Map \mathversion{bold}$F_{\text{cub}} = E(F)$:}
    \begin{align*}
    F_{\text{cub}} = \big( & x_1 - 3x_1^2x_2 + x_7x_9 + x_1^2x_9 + x_1x_3x_7, \\
    & x_2  + 3x_1x_3 + 24x_1x_2^2 - 12x_5x_6 - 36x_2^2x_5 - 12x_8x_{11} \\
    & \qquad - 12x_1x_2x_{11} - 12x_1x_3x_8 + 12x_1x_6x_8 + 12x_3x_5x_8, \\ 
    & x_3 + 8x_2^2 + 6x_1x_2x_3 - 8x_4x_6 - 24x_2^2x_4 - 4x_5x_{10} \\*
    & \qquad - 12x_2x_3x_5 + 8x_2x_5x_6 - 4x_8x_{12} - 4x_1x_2x_{12} \\*
    & \qquad  -28x_2^2x_8 + 8x_3x_4x_8 + 4x_1x_{10}x_8, \\
    & x_4 - x_8^2 - 2x_1x_2x_8, \; x_5 + x_1^2x_2, \; x_6 + x_1x_2x_3 + 3x_2^2, \\
    & x_7 + x_1^2, \; x_8 + x_1x_2, \; x_9 + x_1x_3, \; x_{10} + 3x_2x_3 - 2x_2 x_6, \\
    & x_{11} + x_1x_3 - x_1x_6 - x_3x_5, \; x_{12} + 7x_2^2 - 2x_3x_4 - x_1x_{10} \big)
    \end{align*}
\end{itemize}

\subsection{A Twelve-Dimensional Configuration with Five Cubic Components}
\label{subsec:cubic_12d_five}

Following the implementation of our first refined grid, we investigated whether a unipotent cubic mapping strictly required six cubic components to achieve reduction closure. We discovered that the system could be optimized further by adding $F_1 = 
x_1 - 3x_1^2x_2 - x_1^3x_3$ to the baseline monomial core, resulting in a 12-dimensional coordinate system with only five cubic components (the usefulness of this will become clear in the next subsection). $F_1$ replaces the fourth-degree monomial in the previous constructions within the core grid. 

Under the structural swap of $F_1$ and the quartic monomial in the previous constructions, two cubic monomials remain sufficient to drive the reduction, though the configuration demands one additional quadratic monomial. Because this new quadratic block cannot be linearly combined with any existing core monomials, the system maintains a stable total count of exactly 12 components. 

Furthermore, during the final transition forcing components $F_2$ and $F_3$ from degree 4 down to degree 3, only three remaining fourth-degree terms manifest across the network. By exploiting this low algebraic overhead, we successfully integrate the final sweeping steps directly into the core factorization reduction over the baseline grid, removing the necessity for an independent quotient elimination phase. The complete system parameters, evaluation mappings, and stable integer preimages are extracted directly from our symbolic worksheet \texttt{section4c.mac} as follows:

\begin{itemize}
    \item \textbf{The Initial Extended Map \mathversion{bold}$F$:}
    \begin{align*}
    F = \big( & x_1 - 3x_1^2x_2 - x_1^3x_3, \; 3v^2x_1x_3 + (9v^2-6v-2)x_2, \\
    & v^3x_3 + (6v^2+2v)x_2^2, \; x_4 + x_1^2x_3 + 3x_1x_2, \; x_5 + x_2^3, \\
    & x_6 + x_1^2, \; x_7 + x_1x_3, \; x_8 + x_2^2, \; x_9 + x_1x_5, \\ 
    & x_{10} + x_1x_8, \; x_{11} + x_1x_9, \; x_{12} + x_1x_2 + x_1x_{10} \big)
    \end{align*}
    where $v = 1 + 2x_1x_2$.
    
    \item \textbf{The Stable Geometric Preimages \mathversion{bold}$P$:}
    \begin{align*}
    P = \big\{ & (0, \, 1, \, -8, \, 0, \, -1, \, 0, \, 0, \, -1, \, 0, \, 0, \, 0, \, 0), \\
               & (-1, \, 1, \, 4, \, -1, \, -1, \, -1, \, 4, \, -1, \, -1, \, -1, \, -1, \, 0) \big\}
    \end{align*}
    
    \item \textbf{The Reduction Duality Identity Vector \mathversion{bold}$E$:}
    \begin{align*}
    E = \big( & y_1 + y_6y_7, \; y_2 + 12y_1y_8 + 12y_4y_{10} - 12y_7y_{12}, \\
              & y_3 + 8y_1y_5 - 12y_4y_8 + 8y_4y_9 + 12y_7y_{10} - 8y_7y_{11}, \\
              & y_4, \; y_5, \; y_6, \; y_7, \; y_8, \; y_9, \; y_{10}, \; y_{11}, \; y_{12} \big)
    \end{align*}
    
    \item \textbf{The Resulting Cubic Map \mathversion{bold}$F_{\text{cub}} = E(F)$:}
    \begin{align*}
    F_{\text{cub}} = \big( & x_1 - 3x_1^2x_2 + x_6x_7 + x_1^2x_7 + x_1x_3x_6, \\
    & x_2 + 3x_1x_3 + 24x_1x_2^2 + 12x_1x_8 + 12x_1x_2^2 + 12x_4x_{10} \\
    & \qquad + 36x_1x_2x_{10} + 12x_1x_4x_8 - 12x_7x_{12} - 12x_1x_3x_{12} \\
    & \qquad - 12x_1x_2x_7 - 12x_1x_7x_{10}, \\
    & x_3  + 8x_2^2 + 6x_1x_2x_3 + 8x_1x_5 - 12x_4x_8 - 36x_1x_2x_8 - 12x_2^2x_4 \\
    & \qquad + 8x_4x_9 + 24x_1x_2x_9 + 8x_1x_4x_5 + 12x_7x_{10} + 12x_1x_3x_{10} \\
    & \qquad + 12x_1x_7x_8 - 8x_7x_{11} - 8x_1x_3x_{11} - 8x_1x_7x_9, \\
    & x_4 + x_1^2x_3 + 3x_1x_2, \; x_5 + x_2^3, \; x_6 + x_1^2, \; x_7 + x_1x_3, \; x_8 + x_2^2, \\
    & x_9 + x_1x_5, x_{10} + x_1x_8, \; x_{11} + x_1x_9, \; x_{12} + x_1x_2 + x_1x_{10} \big) 
    \end{align*}
\end{itemize}

\begin{remark}
It is instructive to analyze the algebraic behavior of the first component $F_1$ under the geometric locus. One might initially expect that the condition $F_1 = 0$, which holds at the parameterized coordinate destination, is the primary mechanism that prevents the non-injective collision points from being altered during the transformation. However, a precise tracking of the reduction sequence reveals that $F_1$ is systematically multiplied by the auxiliary components $F_5$ and $F_6$ prior to the elimination phase. Because both $F_5$ and $F_6$ already vanish identically under the kernel initialization, their vanishing is completely sufficient to ensure that the coordinate vectors of the preimages are extended trivially by appending zeroes. Consequently, the condition $F_1 = 0$ is not important for the current elimination track. 

Even if the duality vector $E$ had contained a pure $y_1$-term, such as $y_1^2$, the core injectivity failure would remain entirely unbroken, though it would shift the numerical coordinate positions of the specialized preimages within the extended domain.
\end{remark}

\subsection{The Resulting Stable Unipotent and Homogeneous Counterexamples}
\label{subsec:cubic_stable_mappings}

Following the reduction algorithms detailed above, the primary loop terminates successfully at our three structural milestones of 14, 12 (with six cubic components), and 12 (with five cubic components) variables, respectively. To lift these intermediate cubic mappings to their final stable configurations, we invoke the stabilization algorithms from Section~\ref{sec:unipotent_reduction}. Because forcing uniform degrees across the networks depends linearly on the remaining degree-deficits, the stabilization branches into two completely equivalent pathways across all three frameworks:

\begin{enumerate}
    \item \textbf{The Non-Homogeneous Maps $G$:} Specializing the parameter configurations without enforcing global homogenization collapses the parameter spaces to exact final dimensions evaluated by the sum $M = 2c + q$, where $c$ represents the stabilized cubic variables and $q$ represents the quadratic-lifting parameters.
    \begin{itemize}
        \item For our initial 14-dimensional baseline map ($c=6, q=8$), the dimension collapses to $M = 2 \times 6 + 8 = 20$, yielding an explicit unipotent polynomial map $G: \mathbb{C}^{20} \to \mathbb{C}^{20}$ of degree 3 and a definitive \textbf{40-dimensional non-homogeneous counterexample of degree 4 to the Vanishing Conjecture}.
        \item For our refined 12-dimensional map ($c=6, q=6$), the dimension evaluates to $M = 2 \times 6 + 6 = 18$, yielding an explicit unipotent polynomial map $G: \mathbb{C}^{18} \to \mathbb{C}^{18}$ of degree 3 and a definitive \textbf{36-dimensional non-homogeneous counterexample of degree 4 to the Vanishing Conjecture}.
        \item For our optimal 12-dimensional map ($c=5, q=7$), the dimension collapses to a record minimum of $M = 2 \times 5 + 7 = 17$, yielding an explicit unipotent polynomial map $G: \mathbb{C}^{17} \to \mathbb{C}^{17}$ of degree 3 and a definitive \textbf{34-dimensional non-homogeneous counterexample of degree 4 to the Vanishing Conjecture}.
    \end{itemize}
    
    \item \textbf{The Homogeneous Maps $G_h$:} Alternatively, enforcing strict uniformity of degree 3 under a global scale parameter adds exactly 1 synchronization axis to the core sum, elevating the total architectures to their final homogeneous dimensions:
    \begin{itemize}
        \item For our initial baseline framework, the dimension evaluates to $M = 2 \times 6 + 8 + 1 = 21$, yielding an explicit homogeneous polynomial map $G_h: \mathbb{C}^{21} \to \mathbb{C}^{21}$ and a definitive \textbf{42-dimensional homogeneous counterexample of degree 4 to the Vanishing Conjecture}.
        \item For our refined framework, the dimension collapses to $M = 2 \times 6 + 6 + 1 = 19$, yielding an explicit homogeneous polynomial map $G_h: \mathbb{C}^{19} \to \mathbb{C}^{19}$ and a definitive \textbf{38-dimensional homogeneous counterexample of degree 4 to the Vanishing Conjecture}.
        \item For our optimal framework, the dimension reaches an absolute structural limit of $M = 2 \times 5 + 7 + 1 = 18$, yielding an explicit homogeneous polynomial map $G_h: \mathbb{C}^{18} \to \mathbb{C}^{18}$ and a definitive \textbf{36-dimensional homogeneous counterexample of degree 4 to the Vanishing Conjecture}.
    \end{itemize}
\end{enumerate}

Due to the intensive variable tracking and multi-layered quotient alignments, the full polynomial expansions of these mappings are algebraically too large to be practically displayed here. The complete, uncurated coordinate configurations and automated verification scripts for our reductions are provided in the respective ancillary files \texttt{section4a.mac}, \texttt{section4b.mac} and \texttt{section4c.mac}.

To demonstrate the wider strategic flexibility of our optimal coordinate grid optimization, the dimensional reduction method utilizing five cubic components can also be partially deployed across intermediate degree boundaries. By systematically mapping the degree structures of the component fields, we chart the exact stabilization bounds for maximum degrees $d$ ranging from 7 down to the cubic minimum. The complete automated calculations and parameter chains for these intermediate stabilization tracks are preserved within the verification files \texttt{section2b.mac}, \texttt{section2c.mac}, and \texttt{section2d.mac} respectively. This refined interpolation framework yields the highly optimized multi-degree classification summarized below.
\begin{center}
\begin{tabular}{ccccc}
\hline
\multirow{2}{*}{\textbf{\mathversion{bold}$d$}} & 
\textbf{\mathversion{bold}$n$ such that} & 
\textbf{Degrees of} &
\textbf{\mathversion{bold}$n$ such that} & 
\textbf{\mathversion{bold}$n$ such that} \\
&
\textbf{\mathversion{bold}$\neg \mathrm{JC}(n,d)$} & 
\textbf{Components} &
\textbf{\mathversion{bold}$\neg \mathrm{UJC}(n,d)$} & 
\textbf{\mathversion{bold}$\neg \mathrm{HJC}(n,d)$} \\ 
\hline
7   & 3                                      & 4, 7, 6                     & 14                                      & 15                                      \\
6   & 4                                      & 4, 5, 6, 3                  & 14                                      & 15                                      \\
5   & 5                                      & 4, 5, 5, 3, 2               & 14                                      & 15                                      \\
4   & 8                                      & 4, 4, 4, 3, 3, 2, 2, 2      & 16                                      & 17                                      \\
3   & 12                                     & 3, 3, 3, 3, 3, 2, \dots, 2 & 17                                      & 18                                      \\ \hline
\end{tabular}
\end{center}
This rigorous classification highlights the remarkable efficiency of our stabilization mechanism; even when the maximum degree $d$ is forced down to 3, the resulting homogeneous Keller map collapses into a highly restricted variable space of only 18 dimensions.

In \texttt{section3a.mac}, \texttt{section3b.mac}, \texttt{section3c.mac}, \texttt{section3d.mac}, and \texttt{section3e.mac}, the counterexamples to $\mathrm{UJC}(n,d)$ and $\mathrm{HJC}(n+1,d)$ are reduced to counterexamples to $\mathrm{VC}(2n,d+1)$ and $\mathrm{HVC}(2n+2,d+1)$ respectively, for $d = 7,6,5,4,3$ respectively.

% *************************************************************************************************

\section{The Descent to Dru\.{z}kowski Cubic--Linearity}
\label{sec:druz_linearity}

In this section, we transition from stable unipotent coordinate frameworks to strictly power-linear structures of Dru\.{z}kowski type. To systematically chart the algebraic path toward linearization and prepare the system for subsequent minimization, we introduce a general reduction mechanism that establishes a clear dimensional upper bound for the resulting linear variable space. Crucially, this reduction framework operates universally for any arbitrary system degree $d \ge 2$, allowing higher-degree homogeneous mappings to be driven down to pure linearity without requiring restrictive low-degree assumptions at the core.

\subsection{The Power-Linear Reduction Theorem}
\label{subsec:power_linear_theorem}

To save one stabilization dimension without tweaking the cubic--homogeneous map $G$ in advance, we do not work directly with the vector field $G(X)$ itself. Instead, we embed our counterexample into a modified cubic--linear framework by applying a pairing technique specifically to the perturbed mapping $G(X) + x_n^d e_n$, which isolates the final coordinate axis as an explicit higher-order generator. Although our reduction builds upon structural concepts originally introduced by Gorni and Zampieri \cite{GorniZampieri1997}, the reduction procedure and the proof of Theorem~\ref{thm:druz_main} are entirely self-contained, depending strictly on the inner algebraic properties of our specific coordinate representations.

\begin{theorem} \label{thm:druz_main}
Let $R$ be a commutative Hermite ring and $G: R^n \to R^n$ be a polynomial map, such that $G(X) - X$ is homogeneous of degree $d$, with its last component $G_n(X) - x_n = 0$. Let $H = G(X) - X + x_n^d e_n$. Suppose there exists a polynomial map $\mathit{PL} = (L_1(X)^d, \dots, L_m(X)^d)^T$ of size $m$ such that every component of $H$ can be expressed as a linear combination of the components of $\mathit{PL}$. Let $B \in R^{n \times m}$ be the coefficient matrix defined by $H = B\,\mathit{PL}$. If $B$ is right-invertible, and the last row of $B$ contains an entry in $R^{\times}$, then there exists a Drużkowski-type stabilization of $G$ in exactly dimension $m$.
\end{theorem}

\begin{proof}
Let $C$ be the right-inverse of $B$. We first show that we can extend $B$ and $C$ to square matrices which are the inverse of each other. If $v, w \in R^m$ and $CB v = (I_m - CB) w$, then $CB v = CBCB v = (CB - CBCB) w = (CB - CB) w = 0$, so
$$
R^m = CB \cdot R^m \oplus (I_m - CB) \cdot R^m
$$
Notice that $CB \cdot R^m = CBCB \cdot R^m$ and $R^n = BC \cdot R^n$ are isomorphic via $B$ and $C$. As $R$ is a commutative Hermite ring, the stably free module $(I_m - CB) \cdot R^m$ is isomorphic to $R^{m-n}$. 

Make $C_{\text{ext}}$ from $C$ by adding a basis of $(I_m - CB) \cdot R^m$ as columns to $C$. Since the first $n$ columns of $C_{\text{ext}}$ span $CB \cdot R^m$, it follows that
$$
C_{\text{ext}} \cdot R^m \supseteq CB \cdot R^m + (I_m - CB) \cdot R^m = R^m
$$
so there exists a matrix $B_{\text{ext}}$ such that $C_{\text{ext}} B_{\text{ext}} = I_m$. Hence $v^T C_{\text{ext}} = 0 \Rightarrow v^T = v^T C_{\text{ext}} B_{\text{ext}} = 0$, so $\ker C_{\text{ext}}^T = \{0\}$. If we repeat the extension procedure degenerately on $C_{\text{ext}} B_{\text{ext}} = I_m$ instead of $BC = I_n$, then we obtain that $B_{\text{ext}}$ has a right-inverse $B_{\text{ext}}^{-1}$, so 
$$
B_{\text{ext}}C_{\text{ext}} = B_{\text{ext}}C_{\text{ext}}B_{\text{ext}}B_{\text{ext}}^{-1} = B_{\text{ext}}B_{\text{ext}}^{-1} = I_m
$$
Since $B C_{\text{ext}} = (I_n \,|\, 0_{n \times (m-n)})$ and $\ker C_{\text{ext}}^T = \{0\}$, we infer that $B_{\text{ext}}$ is indeed an extension of $B$. So the first $n$ components of the polynomial map $B_{\text{ext}}\,\mathit{PL}$ are exactly equal to $H$, while the remaining $m-n$ components reside in the polynomial ring $R[X]$.

We now extend the original map $G(X)$ to an $m$-dimensional system $G_{\text{ext}}: R^m \to R^m$ by introducing auxiliary variables $x_{n+1}, \dots, x_m$ and appending suitable linear combinations of the components of $\mathit{PL}$. Specifically, we define:
\[ G_{\text{ext}}(X_{\text{ext}}) = \begin{pmatrix} 0 \\ x_{n+1} \\ \vdots \\ x_m \end{pmatrix} + \begin{pmatrix} G(X) \\ (B_{\text{ext}})_{n+1} \mathit{PL} \\ \vdots \\ (B_{\text{ext}})_m \mathit{PL} \end{pmatrix} \]
where $X_{\text{ext}} = (x_1, \dots, x_m)^T$. 
Multiplying the extended system by $C_{\text{ext}}$ yields:
\[ C_{\text{ext}} G_{\text{ext}}(X_{\text{ext}}) = C_{\text{ext}} X_{\text{ext}} + \mathit{PL} - C_{\text{ext}} e_n x_n^d \]
for some $v \in R^m$.

Suppose that $B_{ni} \in R^{\times}$. To eliminate the problematic $x_n^d$ terms, we construct an invertible matrix $\tilde{C}_{\text{ext}} \in R^{m \times m}$ from $C_{\text{ext}}$ by replacing its $i$-th row with the standard basis vector $e_n^T$. Under this targeted modification, the $i$-th component of the transformed system $\tilde{C}_{\text{ext}} G_{\text{ext}}(X_{\text{ext}})$ becomes identically equal to $x_n$. Because $x_n$ is now isolated within a purely linear coordinate row, we can eliminate the $x_n^d$ terms by way of compositions with elementary maps.

To establish the invertibility of $\tilde{C}_{\text{ext}}$, notice first that the expansion of $\det(\tilde{C}_{\text{ext}})$ along the $i$-th row of $\tilde{C}_{\text{ext}}$ is equal to the $(i,n)$-cofactor of $C_{\text{ext}}$. By Cramer's rule, this cofactor is exactly equal to $\det(C_{\text{ext}}) (B_{\text{ext}})_{ni}$. Since $(B_{\text{ext}})_{ni} = B_{ni} \in R^{\times}$ and $C_{\text{ext}}$ is invertible, $\det(\tilde{C}_{\text{ext}})$ is the product of two units, implying $\det(\tilde{C}_{\text{ext}}) \in R^{\times}$. Again by Cramer's rule, this guarantees that $\tilde{C}_{\text{ext}}^{-1}$ exists over $R$.

Finally, to repair the linear part of the map and bring it into the official Drużkowski format, we apply the linear coordinate transformation $X = \tilde{C}_{\text{ext}}^{-1} Y$, where $Y$ represents the new $m$-dimensional coordinate vector. The substitution yields a map of the exact form $F(Y) = Y + (A Y)^{\times d}$ for a constant matrix $A \in R^{m \times m}$.
\end{proof}

\begin{remark}
To see explicitly how this construction saves a stabilization dimension, consider the algebraic interplay between the final component and the rank condition. Because the original map satisfies $G_n(X) = x_n$, its intrinsic cubic part is identically zero, meaning that without modification, the $n$-th component of the target map would satisfy $H_n = 0$. If the matrix $B$ is to maintain full rank $n$, the fact that $H_n = 0$ mathematically forces the existence of a non-trivial linear relation among the components of the cubic-linear power map $\mathit{PL}$ to prevent a row of zeroes. 

Furthermore, this explanation clarifies why we do not need to pre-emptively alter the system. If $x_n^d$ were already a component of $\mathit{PL}$, one could tweak the cubic-homogeneous map $G$ in advance to avoid generating $x_n^d$ entirely. This would naturally save a dimension as well, without the need to apply the Gorni-Zampieri pairing on $G(X) + x_n^d e_n$ instead of $G(X)$.
\end{remark}

\begin{remark}
Although Theorem \ref{thm:druz_main} establishes the existence of a Drużkowski counterexample in dimension $m$, it does not automatically guarantee that $m$ is the absolute minimal dimension achievable. However, if the resulting Drużkowski map is constructed directly from $G$ and the polynomial map $\mathit{PL}$, one can reasonably expect that all individual components of $\mathit{PL}$ manifest themselves in a transparent, structural manner within the final cubic-linear framework. 

Under the assumption that no linear dependencies exist between the components of $\mathit{PL}$, the condition that $B$ is right-invertible implies that the components of $H$ are linearly independent. While this allows for the explicit generation of $x_n^d$, any further reduction in dimension would necessitate a non-trivial linear intertwining of the independent powers. Consequently, assuming that the essential linear transformations of the components of $\mathit{PL}$ must be preserved within the cubic parts of the Drużkowski map, all $m$ dimensions are strictly required, making the dimension tight for this specific stabilization framework.
\end{remark}

\subsection{A Drużkowski Stabilization of the Cubic Homogeneous Counterexample}

We have applied Theorem \ref{thm:druz_main} on $G = G_h$. The construction of the polynomial map $\mathit{CL}$ has been performed via a two-pass algebraic scanning algorithm:

\begin{enumerate}
    \item \textbf{First Pass (Power Generation):} We isolated the terms of each component of $H$. Terms of the form $x_i^3$ were taken directly, and terms of the form $x_i^2x_j$ were expanded using the identity 
$$
6 x_i^2 x_j = (x_i+x_j)^3 - (x_i-x_j)^3 - 2x_j^3
$$
Purely mixed terms $x_i x_j x_k$ were expanded using 
$$
6 x_i x_j x_k = (x_i + x_j + x_k)^3 - (x_i+x_j)^3 - (x_i+x_k)^3 - (x_j+x_k)^3 + x_i^3 + x_j^3 + x_k^3
$$
instead of the following expansion of $24 x_i x_j x_k$
$$
(x_i + x_j + x_k)^3 + (x_i - x_j - x_k)^3 - (x_i + x_j - x_k)^3 - (x_i - x_j + x_k)^3
$$
in favor of reusing terms.
    \item \textbf{Second Pass (Matrix Assembly):} A coefficient matrix $B$ of size $n \times m$ was directly populated using the exact fractions derived from the algebraic identities, bypassing the computational bottleneck of solving large symbolic systems.
\end{enumerate}

This constructive procedure successfully terminated with a series of compact linear expansions for each of the three cubic--homogeneous mappings under analysis. The exact dimensions $n$ of the homogeneous maps and the resulting numbers $m$ of unique linear powers are summarized below.
\begin{center}
\begin{tabular}{ccc}
\hline
\textbf{Configuration Track} & 
\textbf{Base Dimension (\mathversion{bold}$n$)} & 
\textbf{Linear Powers (\mathversion{bold}$m$)} \\ 
\hline
Baseline Framework          & 21                           & 129                          \\
Refined Grid Track          & 19                           & 107                          \\
Optimal Grid Track          & 18                           & 104                          \\ 
\hline \\
\end{tabular}
\end{center}

The structural and algebraic validity of these reductions were formally verified in \texttt{section4a.mac}, \texttt{section4b.mac}, and \texttt{section4c.mac} respectively. The verification rigorously confirmed three interconnected conditions:

\begin{enumerate}
    \item \textbf{Identity Condition:} The exact matrix-vector identity test was executed via $B \cdot \mathit{CL} \equiv H$. The verification succeeded with a perfect algebraic match across all rows, confirming that the generated powers flawlessly reconstruct the target cubic-homogeneous system.
    \item \textbf{Rank Condition of the System:} The rank of the overarching coefficient matrix $B$ was evaluated, confirming $\operatorname{rank}(B) = n$. This guarantees that the $n$ original components are maximally linearly independent within the $m$-dimensional spanning space.
    \item \textbf{Linear Independence of the Powers:} To guarantee that the dimension cannot be further reduced, we verified the linear independence of the cubic powers themselves. This was achieved by mapping the components of $\mathit{CL}$ onto the standard monomial basis $\mathcal{M}$ of degree 3, consisting of all unique monomials $x_i^3$, $x_i^2 x_j$, and $x_i x_j x_k$ present in the expansions of the components of $H$. We computed a representation matrix $M$ of the coefficients of these monomials for each component of $\mathit{CL}$, and verified that $\operatorname{rank}(M) = m$.
\end{enumerate}

\begin{remark}
It is highly instructive to evaluate these power-linear reductions from a ring-theoretic perspective by generalizing the Dru\.{z}kowski transformation to commutative Hermite rings. While our unipotent expansions in Section~\ref{sec:unipotent_reduction} maintain complete arithmetic purity over the integers, attempting to establish a strictly integer-pure cubic--linear counterexample over the Hermite ring $R = \mathbb{Z}$ encounters a fundamental barrier. Although the standard polarization matrix $B$ inherently introduces fractional coefficients, one can attempt to clear these denominators by scaling the unipotent mapping to the form $G(X) = X + 6(G_h - X)$, which successfully forces all matrix entries of $B$ to lie strictly within $\mathbb{Z}$. However, this modification fails to yield an integer-pure Dru\.{z}kowski reduction because the resulting integral matrix $B$ loses its right-invertibility over $\mathbb{Z}$. 

In our 12-dimensional cubic configurations, this obstruction is driven by the fact that the entire cubic part of the second component $(F_{\text{cub}})_{2}$ is strictly divisible by 12, forcing a factor of 12 into the corresponding minor determinants of $B$. Because 12 is not a unit in $\mathbb{Z}$, the required right-inverse inevitably introduces fractional coefficients. Interestingly, in our baseline 14-dimensional cubic configuration, the cubic part of $(F_{\text{cub}})_{2}$ exhibits two terms that form an exception to this uniform divisibility, namely $-x_1x_2x_{11}$ and $-x_1x_3x_{12}$. However, this structural exception does not bypass the underlying ring obstruction; it merely shifts the divisibility by 12 deeper into the network, as the non-linear parts of both baseline auxiliary equations $(F_{\text{cub}})_{11}$ and $(F_{\text{cub}})_{12}$ remain strictly divisible by 12. Consequently, the descent to pure Drużkowski linearity inherently breaks the integer purity of the counterexamples, necessitating a transition to the field of rational or complex numbers.
\end{remark}

\begin{remark}
\label{rem:druz_integer_pure_construction}
Despite the obstruction preventing the polarization matrix $B$ from achieving full right-invertibility over the integers, it remains structurally possible to guarantee that the final resulting Dru\.{z}kowski cubic--linear mapping itself is defined strictly over $\mathbb{Z}$. A direct, brute-force approach to clearing the denominators from the construction consists of replacing the finalized Dru\.{z}kowski map $F(Y)$ with the scaled configuration $6^{-3} F(6^3 Y)$. Because this substitution uniformly multiplies every non-linear component of the cubic--linear matrix field $AY$ by a factor of $6^2$, it eliminates any lingering fractional elements and successfully restores the arithmetic integer purity of the polynomial endomorphism.

A more subtle and mathematically elegant alternative exploits the scaling degrees of the extended matrix representations. By selecting the modified unipotent map $G(X) = X + 6(G_h - X)$ instead of $G_h$, the associated polarization matrix $B$ becomes a strict matrix over $\mathbb{Z}$. Following the parameterizations established in Section~\ref{sec:unipotent_reduction}, the vector $(0, \, t, \, 0, \, \dots, \, 0, \, 1)$ acts as a parameterized collision point for the deformed mapping $X + t^{-1}(G_h - X)$, with its corresponding preimage coordinates living within the localized ring $\mathbb{Z}[t]$. Specializing this locus by substituting $t = \frac{1}{6}$ yields a valid coordinate collision point for the integer-mapped system with preimages defined strictly over the ring of fractions $\mathbb{Z}[\frac{1}{6}]$. 

Furthermore, since the cubic power $x_n^3$ is explicitly invoked as a component of the power-linear vector, the final row of the polarization matrix $B$ simplifies to the standard basis vector $e_i^T$. This forcing condition dictates that the $i$-th row of the extended matrix $C_{\text{ext}}$ takes the form $e_n^T$, ensuring a direct inverse matching with the extended base $\tilde{C}_{\text{ext}}^{-1} = B_{\text{ext}}$. The final Dru\.{z}kowski transformation maps back to the functional template $F(Y) = Y + (A B_{\text{ext}} Y)^{\times 3}$. While the first $n$ rows of $B$ are defined over $\mathbb{Z}$, the remaining $m-n$ columns of $C$ must span a specific baseline vector space. Because we retain the algebraic freedom to scale these remaining columns of $C$ arbitrarily, the corresponding reciprocal rows of the inverse matrix $B_{\text{ext}}$ scale inversely. Consequently, one can always select a balanced basis scaling such that the complete extended matrix $B_{\text{ext}}$ is driven entirely into the integer ring $\mathbb{Z}$, achieving a fully integral polynomial mapping over the domain. We have explicitly executed and verified this integral power-linear reduction for our 18-dimensional cubic--homogeneous configuration within the accompanying ancillary file \texttt{section4d.mac}.
\end{remark}

\subsection{A Refinement of the Polarization Method}

To reduce the dimension from the initial $m = 104$ boundaries down to exactly $m = 90$, we introduce a refined polarization method that avoids expanding every mixed monomial independently. Instead, we identify coupled blocks of terms within the components $H_3$, $H_{13}$, $H_{14}$, and $H_{15}$ and group them by defining five strategic linear forms, denoted as $y_1, y_2, y_3, y_4, y_5$:
\begin{align*}
y_1 &= 2 x_9 - 3 x_8 & y_2 &= 3 x_{10} - 2 x_{11} & y_3 &= 3 x_2 - x_7 \\
y_4 &= x_2 + x_{10} & y_5 &= 3 x_2 + 6 x_{10} - 4 x_{11} 
\end{align*}

Then two pairs of the $7$ terms of $H_3$ can be coupled, so $H_3$ can be expressed as a sum of $5$ terms:
$$
H_3 = 4 y_1 x_4 x_{18} + 4 y_2 x_7 x_{18} + 8 x_1 x_5 x_{18} + 8 x_2^2  x_18 - x_{15} x_{18}
$$
Both terms of $H_{12}$ can be coupled, so $H_{12}$ can be expressed as one term:
$$
H_{12} = y_4 x_1 x_{18}
$$
Likewise, one pair of the $3$ terms of $H_{13}$ can be coupled, so $H_{13}$ can be expressed as a sum of $2$ terms:
$$
H_{13} = -y_3 x_1^2 + x_1 x_3 x_6
$$
In $H_{14}$, a group of $4$ of the $6$ terms can be coupled, so $H_{14}$ can be expressed as a sum of $3$ terms:
$$
H_{14} = 12 y_3 y_4 x_1 - 12 x_1 x_3 x_{12} + 12 x_1 x_4 x_8
$$
In $H_{15}$, a group of $4$ and another group of $3$ of the $9$ terms can be coupled, so $H_{15}$ can be expressed as a sum of $4$ terms:
$$
H_{15} = 4 y_1 y_3 x_1 + 2 y_5 x_1 x_3 + 8 x_1 x_4 x_5 - 12 x_2^2 x_4
$$

After the above grouping, only $94$ dimensions are necessary, so the grouping reduces the dimension by $10$. But a further reduction by $4$ is accomplished the other way around, namely by two de-groupings, which reduce the dimension by $1$ and $3$ respectively.

The above grouping leaves the term $12 x_4 x_{10} x_{18}$ of $H_2$ as the only term that contains $x_{10}$ directly. The power expansion for $6 x_4 x_{10} x_{18}$ gives $7$ cubic powers of linear forms, of which $4$ contain $x_{10}$. These $4$ cubic powers cannot be reused, so at least $4$ new cubic powers are needed. But we can arrange that only $3$ new cubic powers are needed. For that purpose, we eliminate $x_{10}$ by means of $y_4$:
$$
12 x_4 x_{10} x_{18} = 12 y_4 x_4 x_{18} - 12 x_2 x_4 x_{18}
$$
Notice that $H_3$ contains the term $4 y_1 x_4 x_{18}$, and that $H_{12} = y_4 x_1 x_{18}$, so the cubic powers $(x_4 + x_{18})^3$ and $(y_4 + x_{18})^3$ are ready to be reused. The same holds for $y_4^3$, $x_4^3$ and $x_{18}^3$, so the term $12 y_4 x_4 x_{18}$ only requires $(y_4 + x_4)^3$ and $(y_4 + x_4 + x_{18})^3$ as new cubic powers. Notice again that $H_3$ contains the term $4 y_1 x_4 x_{18}$, that $H_8 = x_2^2 x_{18}$, and that $H_{15}$ contains the term $-12 x_2^2 x_4$, so the cubic powers $(x_4 + x_{18})^3$, $(x_2 + x_{18})^3$ and $(x_2 + x_4)^3$ are ready to be reused. The same holds for $H_{17} = x_2^3$, $x_4^3$, and $x_{18}^3$, so the term $12 x_2 x_4 x_{18}$ only requires $(x_2 + x_4 + x_{18})^3$ as a new cubic power.

The above grouping also leaves $H_{11} = x_1 x_9 x_{18}$ as the only component of $H$ that contains $x_9$ directly. The power expansion for $6 x_1 x_9 x_{18}$ gives $7$ cubic powers of linear forms, of which four contain $x_9$. These four cubic power cannot be reused, so at least four new cubic powers are needed. But we can arrange that only one new cubic power is needed. For that purpose, we eliminate $x_9$ from $H_{11}$ by means of $y_1$:
$$
H_{11} = \tfrac12 y_1 x_1 x_{18} + \tfrac32 x_1 x_8 x_{18}
$$
Since $H_{10} = x_1 x_8 x_{18}$, the second term of $H_{11}$ does not require additional cubic powers, and $(x_1 + x_{18})^3$ is ready to be reused. Notice once again that $H_3$ contains the term $4 y_1 x_4 x_{18}$, and that $H_{15}$ contains the term $4 y_1 y_3 x_1$, so $(y_1 + x_{18})^3$ and $(y_1 + x_1)^3$ are ready to be reused as well. The same holds for $y_1^3$, $x_1^3$ and $x_{18}^3$, so we only need $(y_1 + x_1 + x_{18})^3$ as a new cubic power.

By taking the algebraic displacements and term-coupling principles detailed above, we significantly alter the global dimensional results for the Drużkowski cubic--linear mapping. Through similar algorithmic grouping strategies, the 21-dimensional cubic--homogeneous configuration can be optimized down from its unrefined state to a stable dimension of 117 variables, although its full mathematical layout is omitted here for brevity. Conversely, for our 19-dimensional cubic--homogeneous configuration, a search across the coordinate components and similar term-grouping techniques failed to yield a structural reduction below its established baseline of 107 dimensions.

However, by exploiting the rich multi-term synergetic profile of our optimal 18-dimensional homogeneous map $G_h$, the polarization method drives the system down to an absolute milestone record of exactly 90 variables. The comparative progression of the power-linear reductions across all three evolutionary tracks is systematically summarized below.
\begin{center}
\begin{tabular}{ccc}
\hline
\textbf{Configuration Track} & 
\textbf{Base Dimension (\mathversion{bold}$n$)} & 
\textbf{Linear Powers (\mathversion{bold}$m$)} \\ 
\hline
Baseline Framework          & 21                           & 117                          \\
Refined Grid Track          & 19                           & 107                          \\
Optimal Grid Track          & 18                           & 90                           \\ 
\hline \\
\end{tabular}
\end{center}

The verification matrix $B_y$ confirms with $\text{rank}(B_y) = 90$ that the cubic mapping expands perfectly into this compressed basis, establishing a historic record-low cubic--linear counterexample to the Jacobian Conjecture. The complete automated validation scripts for these systems are preserved across our verification network in \texttt{section4a.mac}, \texttt{section4b.mac}, \texttt{section4c.mac}, and \texttt{section4d.mac} respectively, ensuring immediate numerical reproducibility. In \texttt{section4d.mac}, the stabilization process is tweaked already in the first stage to obtain 
$$
H_{11} = 2x_1 x_9 x_{19} = y_1 x_1 x_{18} + 3 x_1 x_8 x_{18}
$$ 
instead of $H_{11} = x_1 x_9 x_{19} = \tfrac12 y_1 x_1 x_{18} + \tfrac32 x_1 x_8 x_{18}$, to ensure that $B$ is a matrix over $\mathbb{Z}$.

% *************************************************************************************************

\begin{acknowledgement}
The author wishes to express his sincere gratitude to Roy van Rijn and other contributors to Van Rijn's 12-dimensional cubic counterexample and 19-dimensional cubic-homogeneous counterexample. Prior to learning of his independent numerical results, our stabilization framework had settled at a cubic baseline dimension of 14, homogenizing to dimension 21, and a Drużkowski linearization bound of 123 variables. The insight that a 12-dimensional cubic reduction was structurally achievable served as a powerful catalyst. It inspired the deeper algebraic optimizations detailed in this work, which ultimately allowed us to collapse the system to the current historical records of 12 cubic, 18 cubic-homogeneous, and 90 Drużkowski dimensions.

This manuscript was written by Google Chrome AI, under supervision of the author. For each section, the starting point for the AI-model was the corresponding Maxima worksheet. With the exception of section~\ref{sec:druz_linearity}, only small edits have been applied by the author after the AI-model's work. Again with the exception of section~\ref{sec:druz_linearity}, all ideas in the paper are the author's own. The AI-model came with the 7-term formula that we use for $x_i x_j x_k$ and did not mention the 4-term formula. Furthermore, the AI-model spontaneously brought up that we do not need $x_n^3$ in $\mathit{PL}$.
\end{acknowledgement}

\begin{thebibliography}{99}

\bibitem{alpoge2026}
L. Alpöge,
X post, July 20, 2026,
\url{https://x.com/__alpoge__/status/2079028340955197566}.

\bibitem{BCW1982}
H. Bass, E. H. Connell, and D. Wright,
\emph{The Jacobian Conjecture: Reduction of Degree and Formal Expansion of the Inverse},
Bulletin of the American Mathematical Society, Vol. 7, No. 2, pp. 287--330, 1982.

\bibitem{BE2005}
M. de Bondt and A. van den Essen,
\emph{A reduction of the Jacobian conjecture to the symmetric case},
Proceedings of the American Mathematical Society, Vol. 133, No. 8, pp. 2201-2205, 2005.

\bibitem{B2014}
M. de Bondt,
\emph{Symmetric Jacobians},
Central European Journal of Mathematics, Vol. 12, No. 5, pp. 687--710, 2014.

\bibitem{GorniZampieri1997}
Gianluca Gorni and Gaetano Zampieri,
\emph{On cubic-linear polynomial mappings},
Indagationes Mathematicae, Vol. 8, No. 4, pp. 471--492, 1997.

\bibitem{long2026}
X. Long,
\emph{Small Counterexamples to the Gaussian Moments Conjecture},
preprint, arXiv:2607.18186, July 2026.

\bibitem{meng2005}
G. Meng,
\emph{Legendre transform, Hessian conjecture and tree formula},
Applied Mathematics Letters, Vol. 19, No. 6, pp. 503--510, 2006.

\bibitem{mengyang2026}
G. Meng and L. Yang,
\emph{A five-variable counterexample to the Hessian conjecture, and the low-dimensional status of the Jacobian and Hessian conjectures}, preprint, arXiv:2607.22198, July 2026.

\bibitem{Rijn2026}
R. van Rijn. \emph{A twelve-variable degree-three Keller counterexample}, section “Homogeneous
consequence.” In the jacobian-research repository, accessed 15 September
2026. \texttt{https://github.com/
royvanrijn/jacobian-research/blob/main/verified/
TWELVE\_VARIABLE\_DEGREE\_THREE\_KELLER\_COUNTEREXAMPLE.md.}

\bibitem{zhao2007}
W. Zhao,
\emph{Hessian Nilpotent Polynomials and the Jacobian Conjecture},
Transactions of the American Mathematical Society, Vol. 359, No. 1, pp. 249--274, 2007.

\end{thebibliography}

\end{document}